% ============================================================================
% Fichier     : FORN_main.tex
% Titre       : Forensique Numérique
% Auteur      : MINKA MI NGUIDJOÏ Thierry Emmanuel
% Institution : Chaire UNESCO d'Afrique Centrale Accès TIC : ENSP / LIMSI
% Version     : 0.1 -- octobre 2026 (squelette)
% ============================================================================

\documentclass[12pt, twoside, openright, a4paper]{book}

% --- Feuille de style unifiée ---
\usepackage{Style_FORN}

% --- Bibliographie ---
\addbibresource{bibliographie/FORN_references.bib}

% ============================================================================
% MÉTADONNÉES PDF
% ============================================================================
\hypersetup{
    pdftitle={Forensique Numérique},
    pdfauthor={MINKA MI NGUIDJOÏ Thierry Emmanuel},
    pdfsubject={LPI342 · Master Professionnel en Sécurité des Systèmes Informatiques et des Données · Université de Garoua},
    pdfkeywords={forensique numérique, preuve, chaîne de possession, ISO 27037, ISO 27043, acquisition, FORN}
}

\begin{document}

% --------------------------------------------------------------------------
% ÉLÉMENTS PRÉLIMINAIRES (pages romaines)
% --------------------------------------------------------------------------
\frontmatter

    % ==============================================================================
% Fichier : front/00_page_titre.tex
% Titre   : Page de titre - \titreCours
% Auteur  : MINKA MI NGUIDJOÏ Thierry Emmanuel
% Version : 1.1 -- novembre 2026 (charte des cours, modèle de HE2)
% ==============================================================================
% Logos : Université de Garoua à gauche et à droite.
% ==============================================================================

\begin{titlepage}
\thispagestyle{empty}

\begin{tikzpicture}[remember picture, overlay]
    \fill[CouleurPrimaire]
        (current page.north west) rectangle
        ([yshift=-3.6cm] current page.north east);
    \fill[CouleurAccent]
        ([yshift=-3.6cm] current page.north west) rectangle
        ([yshift=-4.0cm] current page.north east);
    \fill[CouleurPrimaire]
        (current page.south west) rectangle
        ([yshift=2.5cm] current page.south east);
    \fill[CouleurAccent]
        ([yshift=2.5cm] current page.south west) rectangle
        ([yshift=2.9cm] current page.south east);
    % En-tête institutionnel, posé dans le bandeau supérieur
    \node[align=center, text=white] at ([yshift=-1.8cm] current page.north) {%
        {\fontsize{11}{13}\selectfont
        \textsc{République du Cameroun}}\\[0.3ex]
        {\color{CouleurAccent!60!white}\rule{0.4\paperwidth}{0.4pt}}\\[0.6ex]
        {\fontsize{10}{12}\selectfont\textbf{\tutelle}}\\[0.2ex]
        {\fontsize{9}{11}\selectfont \faculte{} · \departement}\\[0.2ex]
        {\fontsize{8}{10}\selectfont \unitepedago}};
    % Logos : l'Université de Garoua à gauche et à droite (un seul établissement)
    \node[fill=white, rounded corners=4pt, inner sep=2pt]
        at ([xshift=1.75cm, yshift=-1.8cm] current page.north west)
        {\includegraphics[height=2.6cm]{figures/logos/garoua_logo.png}};
    \node[fill=white, rounded corners=4pt, inner sep=2pt]
        at ([xshift=-1.75cm, yshift=-1.8cm] current page.north east)
        {\includegraphics[height=2.6cm]{figures/logos/garoua_logo.png}};
    % Année et version, posées dans le bandeau inférieur
    \node[align=center] at ([yshift=1.25cm] current page.south) {%
        {\color{white}\bfseries\large Année académique \anneescolaire}\\[0.3ex]
        {\color{CouleurAccent!70!white}\small
        Version 1.0 | Document pédagogique de l'UE \codeUE}};
\end{tikzpicture}

\vspace*{0.6cm}

% --- Titre principal (sur deux lignes) ---
\begin{center}
    {\color{CouleurPrimaire}\rule{0.85\textwidth}{1.2pt}}\\[1.2ex]
    {\fontsize{26}{32}\selectfont\bfseries\color{CouleurPrimaire}
        Forensique\\[0.4ex]
        Numérique}\\[0.8ex]
    {\color{CouleurAccent}\rule{0.85\textwidth}{2pt}}\\[0.5ex]
    {\color{CouleurPrimaire}\rule{0.85\textwidth}{0.5pt}}
\end{center}

\vspace{0.8cm}

\begin{center}
    {\fontsize{13}{16}\selectfont\itshape\color{CouleurSecondaire}
        De la scène à la preuve~:\\
        l'enquête numérique, de l'acquisition au tribunal}
\end{center}

\vspace{0.5cm}

% --- Cartouche sigle et code UE ---
\begin{center}
    \begin{tikzpicture}
        \node[draw=CouleurPrimaire, fill=CouleurPrimaire, rounded corners=6pt,
              inner xsep=14pt, inner ysep=6pt]{%
            \color{white}\bfseries\large \sigleCours{} - UE~\codeUE
        };
    \end{tikzpicture}
\end{center}

\vspace{1.0cm}

% --- Schéma de la discipline : le processus d'investigation numérique ---
\begin{center}
\begin{tikzpicture}[
    f/.style = {draw=CouleurBordDef, fill=CouleurDef, rounded corners=3pt,
        minimum width=2.2cm, minimum height=0.8cm, align=center,
        font=\small\bfseries\color{CouleurPrimaire}},
    arrow/.style = {-{Stealth[length=4pt]}, thick, color=CouleurSecondaire}]
\foreach \t [count=\i] in {Identifier, Préserver, Acquérir, Analyser, Rapporter}
    \node[f] (n\i) at (2.8*\i - 2.8, 0) {\t};
\foreach \i [evaluate=\i as \j using int(\i+1)] in {1,2,3,4}
    \draw[arrow] (n\i) -- (n\j);
\draw[thick, dashed, color=CouleurAccent] (-1.1,-0.75) -- (12.3,-0.75);
\node[font=\footnotesize\itshape, text=CouleurSecondaire] at (5.6,-1.05)
    {chaîne de custody~: intégrité, traçabilité, opposabilité};
\end{tikzpicture}
\end{center}

\vspace{0.6cm}

\begin{center}
    {\large\bfseries\color{CouleurPrimaire}
        \auteurcours}\\[0.5ex]
    {\normalsize\color{CouleurSecondaire}
       Enseignant Chercheur | \tutelle}\\[0.2ex]
    {\small\color{CouleurBordInfo}
        \faculte{} · \departement}\\[0.2ex]
    {\small\color{CouleurBordInfo}
        \href{mailto:maletyon@proton.me}{maletyon@proton.me}}
\end{center}

\end{titlepage}

    \cleardoublepage

    % ==============================================================================
% Fichier : front/01_avant_propos.tex : FORN (LPI342)
% ==============================================================================

\chapter*{Avant-propos}
\addcontentsline{toc}{chapter}{Avant-propos}
\markboth{Avant-propos}{Avant-propos}

Ce livre est le support de l'unité d'enseignement \textbf{LPI342, Forensique
Numérique}, du Master Professionnel en Sécurité des Systèmes Informatiques et
des Données de l'Université de Garoua (Faculté des Sciences, Département de
Mathématiques et Informatiques).

\medskip
La forensique numérique est l'investigation à vocation probatoire~: établir ce
qui s'est passé \emph{et} pouvoir le prouver. Là où le hacking éthique apprend
à simuler l'attaque, et la supervision à la détecter, FORN apprend à
\textbf{enquêter}~: acquérir sans altérer, analyser disque, mémoire, réseau,
mobile et nuage, reconstituer un faisceau, et restituer par un rapport et
devant un tribunal. Le fil conducteur est une enquête unique, l'\textbf{opération
TSANAGA}, menée à Garoua par le cabinet \textbf{NYONG SÉCURITÉ} pour le compte
du parquet, sur les supports saisis chez un ancien employé de l'entreprise
fictive \textbf{LiZbiZ}, soupçonné d'avoir revendu le fichier des clients.

\medskip
Chaque notion est ancrée à la fois dans les standards internationaux (ISO/IEC
27037, 27042 et 27043, NIST SP 800-86 et 800-101, RFC 3227, guide ACPO) et
dans le droit camerounais et africain (lois n°~2010/012 et 2024/017,
Conventions de Malabo et de Budapest). Un fil théorique parcourt le livre~: le
\textbf{Trilemme CRO} (Confidentialité, Fiabilité, Opposabilité) et la
\textbf{chaîne de possession}, par lesquels se juge toute trace.

\medskip
Les cas sont fictifs~: les noms, les dates et les montants le sont~; les
mécanismes techniques et les références normatives sont exacts. Le livre
s'accompagne de l'environnement \textbf{NDAP SUKULU}, qui fait pratiquer
chaque chapitre.

\bigskip
\noindent \IfFileExists{../../Logos/signature_avant_propos.png}{\includegraphics[height=1.6cm]{../../Logos/signature_avant_propos.png}}{\textit{MINKA MI NGUIDJOÏ Thierry Emmanuel}}

    \cleardoublepage

    % ==============================================================================
% Fichier : front/02_mode_emploi.tex : FORN (LPI342)
% ==============================================================================

\chapter*{Mode d'emploi}
\addcontentsline{toc}{chapter}{Mode d'emploi}
\markboth{Mode d'emploi}{Mode d'emploi}

\section*{À qui s'adresse ce cours}

Aux étudiants du Master Professionnel en Sécurité des Systèmes Informatiques
et des Données, et aux praticiens qui exercent ou se destinent aux métiers
d'\textbf{analyste forensique}, d'\textbf{expert judiciaire},
d'\textbf{enquêteur en cybercriminalité} et de \textbf{juriste TIC}. Des
notions de réseaux, de systèmes d'exploitation et de droit numérique sont
utiles.

\section*{Comment lire ce livre}

Le livre comprend un prologue et six parties. Chaque chapitre suit le même
plan~: épigraphe, niveau et prérequis, objectifs, question directrice, sections,
« à retenir », dix exercices gradués, transition. Les niveaux des exercices
sont signalés par des étoiles~: \niveaubase{} fondamental, \niveaumoyen{}
intermédiaire, \niveauexpert{} expert.

\section*{Les encadrés}

\begin{itemize}
    \item \textbf{Scénario opérationnel}~: un épisode de l'opération TSANAGA,
          avec \emph{constatation} (faits) et \emph{analyse} (interprétation,
          niveau de confiance).
    \item \textbf{Artefact}~: la fiche d'une pièce (nature, source, pertinence,
          fiabilité, opposabilité).
    \item \textbf{Chaîne de custody}~: une procédure de préservation de la
          preuve.
    \item \textbf{Analyse CRO}~: l'arbitrage Confidentialité, Fiabilité,
          Opposabilité.
    \item \textbf{Application métier}~: ce que le chapitre change pour un
          métier.
    \item \textbf{Cadre juridique}~: le texte de droit applicable.
\end{itemize}

\section*{Constatation et analyse}

Le cours distingue systématiquement ce que l'on \textbf{constate} (un fait
vérifiable) de ce que l'on en \textbf{déduit} (une interprétation, assortie
d'un niveau de confiance). C'est la discipline première de l'enquête.

\section*{NDAP SUKULU et évaluation}

L'environnement compagnon NDAP SUKULU propose, pour chaque chapitre, une fiche
d'auto-évaluation et, pour les chapitres pratiques, un lab sur pièces. Fiches
et labs se font hors des 30~heures de cours. Le code du cours est
\textbf{\codeUE}.

    \cleardoublepage

    % ==============================================================================
% Fichier : front/03_ndap_sukulu.tex : FORN (LPI342)
% ==============================================================================

\chapter*{NDAP SUKULU, l'environnement compagnon}
\addcontentsline{toc}{chapter}{NDAP SUKULU, l'environnement compagnon}
\markboth{NDAP SUKULU}{NDAP SUKULU}

\begin{center}
    \includegraphics[width=0.52\textwidth]{figures/logos/logo_ndap_sukulu.png}
\end{center}

\medskip
\textit{Ndap sukulu} signifie « la salle de classe » en bassa, et
\textit{Malet Yon}, « ton enseignant ». \textbf{NDAP SUKULU} est
l'environnement compagnon du cours~: pour chaque chapitre, une fiche
d'auto-évaluation de 40~questions, corrigée automatiquement~; pour les
chapitres pratiques, un lab sur pièces (fiches de scellé, images décrites,
journaux, flux exportés) dans un cadre entièrement fermé~; et un tableau de
bord de progression.

\medskip
On y accède par le code de cours \textbf{\codeUE}. Les fiches et les labs se
font hors des 30~heures de cours. Les pièces des labs sont fictives et
établies pour l'enseignement.

% ==============================================================================
% Fichier : front/03b_labs.tex : FORN (LPI342)
% ==============================================================================

\section*{Les dix labs de NDAP SUKULU}

\begin{tableaucours}{L{2.4cm}L{8.4cm}L{1.1cm}}{%
    \tableauHeader{Lab} &
    \tableauHeader{Sujet} &
    \tableauHeader{Ch.}
}
FORN-LAB-01 & Chaîne de possession et scellés & 5 \\
FORN-LAB-02 & Duplication et vérification par hachage & 8 \\
FORN-LAB-03 & Forensique disque~: lire un système de fichiers & 10 \\
FORN-LAB-04 & Fichiers supprimés et signatures & 11 \\
FORN-LAB-05 & Analyse d'une capture mémoire & 12 \\
FORN-LAB-06 & Forensique réseau~: suivre une exfiltration (flux exportés) & 13 \\
FORN-LAB-07 & Forensique mobile~: lire des artefacts & 14 \\
FORN-LAB-08 & Évaluer la recevabilité d'une preuve & 6 \\
FORN-LAB-09 & Reconstitution chronologique multi-domaines & 17 \\
FORN-LAB-10 & Structurer un rapport forensique & 18 \\
\end{tableaucours}

\medskip
Chaque lab dure environ une heure et se corrige automatiquement. Les pièces
se trouvent dans le dossier \texttt{preuves} (lecture seule) et une copie de
travail dans \texttt{travail}.


    \cleardoublepage

    \tableofcontents
    \cleardoublepage

    \listoffigures
    \listoftables
    \cleardoublepage

% --------------------------------------------------------------------------
% CORPS DU COURS (pages arabes)
% Numérotation des chapitres CONTINUE sur tout le cours : le prologue est le
% chapitre 0, puis C01 à C21, un numéro par fichier.
% --------------------------------------------------------------------------
\mainmatter

    \part*{Prologue Opérationnel}
    \addcontentsline{toc}{part}{Prologue Opérationnel}
        \setcounter{chapter}{-1}
        \IfFileExists{partie0/P00_scene_numerique.tex}{% ============================================================================
% Fichier  : partie0/P00_scene_numerique.tex : FORN (LPI342)
% ============================================================================

\chapter{Prologue : la Scène Numérique}
\label{chap:prologue}
\chaptermark{Prologue}

\epigraph{%
    « Tout contact laisse une trace. »%
}{Edmond Locard, 1920}

\niveauChapitre{Introductif}{Aucun prérequis. Les notions sont reprises en
détail dans les chapitres.}

\begin{boxobjectifs}
\begin{itemize}
    \item Découvrir la posture de l'enquêteur numérique.
    \item Distinguer l'enquête de l'attaque, la trace du jugement.
    \item Découvrir le fil rouge LiZbiZ et l'opération TSANAGA.
    \item Situer le Trilemme CRO et la chaîne de possession.
    \item Formuler les questions auxquelles le cours répond.
\end{itemize}
\end{boxobjectifs}

\bigskip

% ============================================================================
\section*{Un ordinateur et une clé USB, scellés dans un sac}
% ============================================================================

Après le rançongiciel de mars, LiZbiZ découvre qu'un ancien employé de son
site de Garoua, M.~NDONGO, aurait, des mois plus tôt, copié puis revendu le
fichier des clients. Le 20~avril 2026, sur réquisition du parquet, un
ordinateur portable et une clé USB sont saisis à son domicile et remis, sous
scellé, au cabinet \sigle{NYONG SÉCURITÉ}, désigné comme expert.

\medskip
Sur ces supports se trouvent, peut-être, la preuve de ce qui s'est passé.
Mais une preuve mal recueillie, mal conservée, ne vaut rien devant un
tribunal\footnote{Cas fictif établi pour le cours. Les noms, les dates et les
montants sont fictifs~; les mécanismes techniques et les références
normatives sont exacts.}. Toute l'enquête tient dans une tension~: établir la
vérité \emph{et} pouvoir la prouver.

% ============================================================================
\section{L'Enquêteur n'est pas un Pirate}
\label{sec:p00-posture}
% ============================================================================

\begin{boxdef}
\termecle{Forensique numérique}~: recherche méthodique, documentée et
recevable des traces laissées par un événement sur les supports numériques,
afin d'établir ce qui s'est passé, et de pouvoir le démontrer devant un
tiers. Elle se distingue du piratage par le cadre (autorisation, droit) et la
posture (neutralité, traçabilité).
\end{boxdef}

\begin{tableaucours}{L{4.0cm}L{4.2cm}L{4.0cm}}{%
    \tableauHeader{} &
    \tableauHeader{Le pirate} &
    \tableauHeader{L'enquêteur forensique}
}
But & Accéder, nuire, profiter & Établir et prouver \\
Cadre & Illégal & Autorisé, encadré \\
Trace & Les efface & Les préserve \\
Posture & Partisane & Neutre, impartiale \\
\end{tableaucours}
\vspace{-0.8em}
\begin{center}{\small\itshape Tableau~0.1. Deux usages des mêmes compétences}\end{center}

% ============================================================================
\section{Le Fil Rouge : LiZbiZ et l'Opération TSANAGA}
\label{sec:p00-filrouge}
% ============================================================================

Le cours suit une même enquête, l'\sigle{opération TSANAGA}, menée à Garoua
par le cabinet NYONG SÉCURITÉ pour le compte du parquet, sur les supports
saisis chez M.~NDONGO.

\begin{tableaucours}{L{3.4cm}L{8.8cm}}{%
    \tableauHeader{Acteur} &
    \tableauHeader{Rôle}
}
M.~BIKORO & Expert forensique, cabinet NYONG SÉCURITÉ \\
Mme~NKENG & Analyste forensique \\
Me~TSALA & Avocat~; veille à la recevabilité \\
Le parquet de Garoua & Autorité requérante \\
M.~NDONGO & Mis en cause (ancien employé de LiZbiZ) \\
\end{tableaucours}
\vspace{-0.8em}
\begin{center}{\small\itshape Tableau~0.2. Les acteurs de l'opération TSANAGA}\end{center}

% ============================================================================
\section{Deux Fils : CRO et Chaîne de Possession}
\label{sec:p00-cro}
% ============================================================================

\begin{boxCRO}
\textbf{Le Trilemme CRO.} \termecle{Confidentialité}~: ne saisir et n'exposer
que le nécessaire. \termecle{Fiabilité}~: la trace est-elle fidèle, intègre,
non altérée~? \termecle{Opposabilité}~: pourra-t-on s'en prévaloir devant le
juge~? Chaque acte d'enquête s'arbitre sur ces trois dimensions, et se
\emph{documente}.
\end{boxCRO}

\begin{boxprocedure}
\textbf{La chaîne de possession, dès la saisie}
\begin{enumerate}
    \item Identifier chaque support (numéro de scellé, description).
    \item Consigner qui l'a, depuis quand, pour quoi.
    \item Ne jamais rompre la continuité de la détention.
    \item Calculer les empreintes avant toute analyse.
    \item Joindre la chaîne au dossier.
\end{enumerate}
\end{boxprocedure}

% ============================================================================
\section{Les Questions du Cours}
\label{sec:p00-questions}
% ============================================================================

\begin{boxinfo}
\textbf{Ce que le cours apprend à répondre.}
\begin{itemize}
    \item Qu'est-ce qu'une enquête numérique, et qu'est-ce qui la distingue
          d'un piratage (Partie~I)~?
    \item Qu'est-ce qui rend une preuve recevable (Partie~II)~?
    \item Comment acquérir sans altérer (Partie~III)~?
    \item Comment analyser disque, mémoire, réseau, mobile, nuage
          (Partie~IV)~?
    \item Comment déjouer l'anti-forensique et reconstituer (Partie~V)~?
    \item Comment restituer, par un rapport et devant un tribunal
          (Partie~VI)~?
\end{itemize}
\end{boxinfo}

\separateur

\begin{boxresume}
\textbf{Ce qu'il faut retenir de ce prologue~:}
\begin{itemize}
    \item \textbf{Tout contact laisse une trace}~; l'enquête la cherche et la
          préserve.
    \item \textbf{L'enquêteur n'est pas un pirate}~: cadre et posture le
          distinguent.
    \item \textbf{Le fil rouge}~: LiZbiZ et l'opération TSANAGA à Garoua.
    \item \textbf{CRO et chaîne de possession} encadrent chaque acte.
    \item \textbf{Établir et prouver}~: la double exigence de l'enquête.
\end{itemize}
\end{boxresume}

\begin{boxinfo}
\textbf{Transition.} Commençons par définir la discipline et sa posture~:
c'est l'objet du Chapitre~\ref{chap:forensique}.
\end{boxinfo}
}{}

    \part{Fondements de la Forensique Numérique}
    \label{part:fondements}

        \IfFileExists{partie1/P01_C01_forensique_numerique.tex}{% ============================================================================
% Fichier  : partie1/P01_C01_forensique_numerique.tex : FORN (LPI342)
% Sources : ISO/IEC 27037, 27042, 27043 ; Locard ; ACPO. Opération TSANAGA.
% ============================================================================

\chapter{Qu'est-ce que la Forensique Numérique}
\label{chap:forensique}
\chaptermark{La forensique numérique}

\epigraph{%
    « La trace numérique est fragile~; c'est le soin de l'enquêteur\\
    qui la rend solide. »%
}{Principe du cours}

\niveauChapitre{Fondamental}{Le prologue~; notions de systèmes et de réseaux.}

\begin{boxobjectifs}
\begin{itemize}
    \item Définir la forensique numérique et ses termes voisins.
    \item Situer son histoire et ses figures fondatrices.
    \item Distinguer ses domaines (disque, mémoire, réseau, mobile, nuage).
    \item Comprendre son double ancrage, scientifique et juridique.
    \item Situer les référentiels (ISO/IEC 27037, 27042, 27043, ACPO).
\end{itemize}
\end{boxobjectifs}

\bigskip

% ============================================================================
\section*{Question directrice : qu'est-ce qui distingue l'enquêteur numérique d'un informaticien ordinaire ?}
% ============================================================================

Un informaticien sait lire un disque~; un enquêteur forensique le lit
\emph{sans l'altérer}, documente chaque geste, et sépare ce qu'il constate de
ce qu'il en déduit. Ce chapitre définit la discipline, son histoire et ses
domaines.

\medskip
\textbf{Plan du chapitre.}
\begin{enumerate}
    \item Définir la discipline (section~\ref{sec:c01-definir}).
    \item Histoire et figures (section~\ref{sec:c01-histoire}).
    \item Les domaines (section~\ref{sec:c01-domaines}).
    \item Le double ancrage (section~\ref{sec:c01-ancrage}).
\end{enumerate}

\separateur

% ============================================================================
\section{Définir la Discipline}
\label{sec:c01-definir}
% ============================================================================

\begin{boxdef}
\termecle{Investigation numérique}~: démarche générale de collecte, d'analyse
et d'interprétation d'éléments numériques pour éclairer une situation, dans
tout contexte (judiciaire, administratif, privé). La \termecle{forensique
numérique} (ou criminalistique numérique) en est le sous-ensemble à vocation
probatoire, soumis à la rigueur de la preuve.
\end{boxdef}

\begin{tableaucours}{L{4.4cm}L{7.8cm}}{%
    \tableauHeader{Terme} &
    \tableauHeader{Sens}
}
Investigation numérique & Démarche générale, tous contextes \\
Forensique (digital forensics) & Cadre probatoire, rigueur de la preuve \\
Criminalistique numérique & Insiste sur la dimension scientifique \\
Recherche de la preuve & Finalité juridique \\
\end{tableaucours}
\vspace{-0.8em}
\begin{center}{\small\itshape Tableau~1.1. Des termes voisins}\end{center}

% ============================================================================
\section{Histoire et Figures}
\label{sec:c01-histoire}
% ============================================================================

La forensique numérique hérite des sciences de la trace. \termecle{Edmond
Locard} (France, 1910) en énonce le principe fondateur~: « tout contact
laisse une trace »~\cite{locard1920}. Appliqué au numérique, il signifie que
toute action sur un système y laisse des marques, pour qui sait les lire et
les préserver.

% ============================================================================
\section{Les Domaines}
\label{sec:c01-domaines}
% ============================================================================

\begin{tableaucours}{L{3.4cm}L{8.8cm}}{%
    \tableauHeader{Domaine} &
    \tableauHeader{Objet}
}
Disque & Systèmes de fichiers, fichiers supprimés, artefacts \\
Mémoire & État vivant du système, secrets, code \\
Réseau & Flux, exfiltration, pilotage, chronologie \\
Mobile & Téléphones~: messages, applications, géolocalisation \\
Nuage et applications & Données hébergées, journaux de services \\
\end{tableaucours}
\vspace{-0.8em}
\begin{center}{\small\itshape Tableau~1.2. Les domaines de la forensique numérique}\end{center}

% ============================================================================
\section{Le Double Ancrage}
\label{sec:c01-ancrage}
% ============================================================================

La forensique numérique s'appuie sur des référentiels qui lui donnent sa
rigueur~: l'\norme{ISO/IEC 27037} (acquisition et préservation),
l'\norme{ISO/IEC 27042} (analyse et interprétation), l'\norme{ISO/IEC 27043}
(principes de l'investigation), et le guide \sigle{ACPO}. Elle s'ancre aussi
dans le droit (Partie~II).

\begin{boxscenario}{Opération TSANAGA. Une discipline, pas une intuition}
\textbf{Constatation.} Devant l'ordinateur saisi, M.~BIKORO ne l'allume pas~:
il le photographie scellé, note son état, et prépare l'acquisition selon
l'ISO/IEC 27037.

\medskip
\textbf{Analyse} (confiance élevée). Le réflexe « ouvrir pour voir » aurait
altéré la preuve. La discipline impose la méthode avant la curiosité~: c'est
ce qui sépare l'enquête de la manipulation.
\end{boxscenario}

\begin{boxmetier}{Analyste forensique}
\textbf{La méthode avant l'outil.} L'analyste connaît les outils, mais c'est
sa méthode (préserver, documenter, séparer fait et interprétation) qui donne
à son travail une valeur probante.
\end{boxmetier}

\separateur

\begin{boxresume}
\textbf{Ce qu'il faut retenir de ce chapitre~:}
\begin{itemize}
    \item \textbf{La forensique} est l'investigation numérique à vocation
          probatoire.
    \item \textbf{Locard}~: « tout contact laisse une trace ».
    \item \textbf{Cinq domaines}~: disque, mémoire, réseau, mobile, nuage.
    \item \textbf{Double ancrage}~: scientifique (ISO 27037/42/43, ACPO) et
          juridique.
    \item \textbf{La méthode} distingue l'enquêteur de l'informaticien
          ordinaire.
\end{itemize}
\end{boxresume}

\begin{boxinfo}
\textbf{Points de contrôle de l'élève.} Je sais~: définir la forensique et
ses termes voisins~; citer Locard~; nommer les domaines~; situer les
référentiels~; dire ce qui distingue l'enquêteur.
\end{boxinfo}

\bigskip

\begin{boxexo}[title={\ding{45}\quad Exercice~1.1. Termes \niveaubase}]
Distinguez investigation numérique et forensique numérique.
\end{boxexo}

\begin{boxexo}[title={\ding{45}\quad Exercice~1.2. Locard \niveaubase}]
Énoncez le principe de Locard et son sens pour le numérique.
\end{boxexo}

\begin{boxexo}[title={\ding{45}\quad Exercice~1.3. Domaines \niveaubase}]
Citez les cinq domaines de la forensique numérique.
\end{boxexo}

\begin{boxexo}[title={\ding{45}\quad Exercice~1.4. Pirate ou enquêteur \niveaumoyen}]
Sur un même disque, qu'est-ce qui distingue le geste du pirate de celui de
l'enquêteur~?
\end{boxexo}

\begin{boxexo}[title={\ding{45}\quad Exercice~1.5. Référentiels \niveaumoyen}]
Associez à chaque norme son objet~: 27037, 27042, 27043.
\end{boxexo}

\begin{boxexo}[title={\ding{45}\quad Exercice~1.6. TSANAGA \niveaumoyen}]
Pourquoi M.~BIKORO n'allume-t-il pas l'ordinateur saisi~?
\end{boxexo}

\begin{boxexo}[title={\ding{45}\quad Exercice~1.7. Double ancrage \niveaumoyen}]
Expliquez le double ancrage, scientifique et juridique, de la discipline.
\end{boxexo}

\begin{boxexo}[title={\ding{45}\quad Exercice~1.8. Criminalistique \niveaumoyen}]
Pourquoi parle-t-on de « criminalistique » numérique~?
\end{boxexo}

\begin{boxexo}[title={\ding{45}\quad Exercice~1.9. Posture \niveauexpert}]
Rédigez une courte charte de la posture de l'enquêteur numérique.
\end{boxexo}

\begin{boxexo}[title={\ding{45}\quad Exercice~1.10. Panorama \niveauexpert}]
Pour l'opération TSANAGA, quels domaines seront probablement mobilisés, et
pourquoi~?
\end{boxexo}

\bigskip

\begin{boxinfo}
\textbf{Pour aller plus loin.}
\begin{itemize}
    \item ISO/IEC 27037~\cite{iso27037}, 27043~\cite{iso27043}~;
          ACPO~\cite{acpo2012}~; Locard~\cite{locard1920}.
\end{itemize}
\end{boxinfo}

\begin{boxinfo}
\textbf{Transition.} Toute enquête suit une démarche~: le
Chapitre~\ref{chap:demarche} en présente les phases.
\end{boxinfo}
}{}
        \IfFileExists{partie1/P01_C02_demarche_enquete.tex}{% ============================================================================
% Fichier  : partie1/P01_C02_demarche_enquete.tex : FORN (LPI342)
% Sources : ISO/IEC 27043 ; NIST SP 800-86 ; ACPO. Opération TSANAGA.
% ============================================================================

\chapter{La Démarche d'Enquête}
\label{chap:demarche}
\chaptermark{La démarche d'enquête}

\epigraph{%
    « Une enquête n'est pas une suite d'outils~;\\
    c'est une suite de décisions documentées. »%
}{Principe du cours}

\niveauChapitre{Fondamental}{Chapitre~\ref{chap:forensique}.}

\begin{boxobjectifs}
\begin{itemize}
    \item Décrire les phases d'une enquête numérique.
    \item Comprendre la planification et l'autorisation.
    \item Situer identification, acquisition, analyse, restitution.
    \item Distinguer constatation et interprétation.
    \item Documenter chaque phase.
\end{itemize}
\end{boxobjectifs}

\bigskip

% ============================================================================
\section*{Question directrice : par où commence une enquête sur les supports de M. NDONGO ?}
% ============================================================================

Devant les scellés, la tentation est d'« aller voir ». L'enquête impose un
ordre~: planifier, identifier, acquérir, analyser, restituer. Ce chapitre
pose cette démarche.

\medskip
\textbf{Plan du chapitre.}
\begin{enumerate}
    \item Les phases (section~\ref{sec:c02-phases}).
    \item Planifier et autoriser (section~\ref{sec:c02-planifier}).
    \item Constatation et interprétation (section~\ref{sec:c02-constat}).
    \item Documenter (section~\ref{sec:c02-documenter}).
\end{enumerate}

\separateur

% ============================================================================
\section{Les Phases}
\label{sec:c02-phases}
% ============================================================================

\begin{tableaucours}{L{3.4cm}L{8.8cm}}{%
    \tableauHeader{Phase} &
    \tableauHeader{Ce qu'on y fait}
}
Planifier & Cadre, autorisation, périmètre, outils \\
Identifier & Repérer et préserver les sources de preuve \\
Acquérir & Dupliquer, calculer les empreintes, tracer \\
Analyser & Examiner la copie, interpréter, recouper \\
Restituer & Rapport, présentation, conservation \\
\end{tableaucours}
\vspace{-0.8em}
\begin{center}{\small\itshape Tableau~2.1. Les phases d'une enquête~\cite{iso27043,nist80086}}\end{center}

% ============================================================================
\section{Planifier et Autoriser}
\label{sec:c02-planifier}
% ============================================================================

\begin{boxjuris}
Une enquête forensique suppose un fondement~: réquisition judiciaire, mandat,
ou autorisation écrite du responsable pour une enquête interne. Au Cameroun,
accéder à un support sans titre expose aux sanctions de la
\loicam{loi n°~2010/012} (Chapitre~\ref{chap:droit}). On vérifie son titre
avant d'agir.
\end{boxjuris}

% ============================================================================
\section{Constatation et Interprétation}
\label{sec:c02-constat}
% ============================================================================

\begin{boxalert}
\textbf{Violation typique.} Écrire « M.~NDONGO a volé les données » dès la
première trace.

\smallskip
\textbf{Règle~:} on constate (« le fichier clients a été copié sur une clé le
3~janvier »), puis on interprète séparément, avec un niveau de confiance. Le
rapport sépare toujours le fait établi de l'hypothèse.
\end{boxalert}

% ============================================================================
\section{Documenter}
\label{sec:c02-documenter}
% ============================================================================

\begin{boxscenario}{Opération TSANAGA. Le journal d'enquête}
\textbf{Constatation.} Dès la réception des scellés, Mme~NKENG ouvre un
journal d'enquête~: chaque acte y est horodaté, décrit, signé.

\medskip
\textbf{Analyse} (confiance élevée). Le journal n'est pas une formalité~:
c'est lui qui rendra l'enquête reproductible et opposable. Une analyse non
documentée, même exacte, sera contestée au tribunal.
\end{boxscenario}

\begin{boxmetier}{Enquêteur en cybercriminalité}
\textbf{Suivre la démarche sous pression.} Même pressé par le parquet,
l'enquêteur ne saute pas de phase~: planifier et identifier avant d'acquérir,
acquérir avant d'analyser. L'ordre protège la preuve.
\end{boxmetier}

\separateur

\begin{boxresume}
\textbf{Ce qu'il faut retenir de ce chapitre~:}
\begin{itemize}
    \item \textbf{Cinq phases}~: planifier, identifier, acquérir, analyser,
          restituer.
    \item \textbf{On vérifie son titre} avant d'agir.
    \item \textbf{On constate, puis on interprète}, séparément.
    \item \textbf{On documente} chaque acte dans un journal d'enquête.
    \item \textbf{L'ordre des phases} protège la preuve.
\end{itemize}
\end{boxresume}

\begin{boxinfo}
\textbf{Points de contrôle de l'élève.} Je sais~: énoncer les phases~;
expliquer le fondement d'une enquête~; séparer constatation et
interprétation~; tenir un journal d'enquête.
\end{boxinfo}

\bigskip

\begin{boxexo}[title={\ding{45}\quad Exercice~2.1. Les phases \niveaubase}]
Citez les cinq phases d'une enquête numérique dans l'ordre.
\end{boxexo}

\begin{boxexo}[title={\ding{45}\quad Exercice~2.2. Le titre \niveaubase}]
Quel fondement autorise l'enquête sur les supports de M.~NDONGO~?
\end{boxexo}

\begin{boxexo}[title={\ding{45}\quad Exercice~2.3. Constat ou interprétation \niveaubase}]
Classez~: « le fichier a été copié le 3~janvier »~; « M.~NDONGO est
coupable ».
\end{boxexo}

\begin{boxexo}[title={\ding{45}\quad Exercice~2.4. Sauter une phase \niveaumoyen}]
Que risque-t-on à analyser avant d'acquérir~?
\end{boxexo}

\begin{boxexo}[title={\ding{45}\quad Exercice~2.5. Journal \niveaumoyen}]
Que consigne-t-on dans un journal d'enquête~?
\end{boxexo}

\begin{boxexo}[title={\ding{45}\quad Exercice~2.6. Confiance \niveaumoyen}]
Pourquoi assortir une interprétation d'un niveau de confiance~?
\end{boxexo}

\begin{boxexo}[title={\ding{45}\quad Exercice~2.7. Pression \niveaumoyen}]
Le parquet veut un résultat vite. Comment tenir la démarche~?
\end{boxexo}

\begin{boxexo}[title={\ding{45}\quad Exercice~2.8. Reproductible \niveaumoyen}]
Pourquoi une enquête doit-elle être reproductible~?
\end{boxexo}

\begin{boxexo}[title={\ding{45}\quad Exercice~2.9. Plan d'enquête \niveauexpert}]
Rédigez le plan d'enquête de l'opération TSANAGA, phase par phase.
\end{boxexo}

\begin{boxexo}[title={\ding{45}\quad Exercice~2.10. Modèle de journal \niveauexpert}]
Concevez le modèle de journal d'enquête du cabinet NYONG SÉCURITÉ.
\end{boxexo}

\bigskip

\begin{boxinfo}
\textbf{Pour aller plus loin.}
\begin{itemize}
    \item ISO/IEC 27043~\cite{iso27043}~; NIST SP~800-86~\cite{nist80086}~;
          ACPO~\cite{acpo2012}.
\end{itemize}
\end{boxinfo}

\begin{boxinfo}
\textbf{Transition.} La démarche n'a de valeur que portée par une éthique~:
le Chapitre~\ref{chap:ethique} traite de la neutralité et de l'impartialité.
\end{boxinfo}
}{}
        \IfFileExists{partie1/P01_C03_ethique_neutralite.tex}{% ============================================================================
% Fichier  : partie1/P01_C03_ethique_neutralite.tex : FORN (LPI342)
% Sources : ISO/IEC 27037 (impartialité) ; ACPO ; déontologie de l'expertise.
% ============================================================================

\chapter{Éthique, Neutralité, Impartialité}
\label{chap:ethique}
\chaptermark{Éthique et impartialité}

\epigraph{%
    « L'enquêteur ne cherche pas à avoir raison~;\\
    il cherche ce qui s'est passé. »%
}{Principe du cours}

\niveauChapitre{Fondamental}{Chapitre~\ref{chap:demarche}.}

\begin{boxobjectifs}
\begin{itemize}
    \item Comprendre la neutralité et l'impartialité de l'enquêteur.
    \item Reconnaître les biais qui menacent l'enquête.
    \item Respecter la confidentialité et la proportionnalité.
    \item Gérer les conflits d'intérêts.
    \item Rapporter à charge et à décharge.
\end{itemize}
\end{boxobjectifs}

\bigskip

% ============================================================================
\section*{Question directrice : l'expert doit-il aider à condamner M. NDONGO ?}
% ============================================================================

L'expert n'est ni du côté de l'accusation ni de la défense~: il éclaire. Son
devoir est la vérité, rapportée à charge comme à décharge. Ce chapitre pose
l'éthique de l'enquête.

\medskip
\textbf{Plan du chapitre.}
\begin{enumerate}
    \item Neutralité et impartialité (section~\ref{sec:c03-neutralite}).
    \item Les biais (section~\ref{sec:c03-biais}).
    \item Confidentialité et proportionnalité (section~\ref{sec:c03-confid}).
    \item Conflits d'intérêts (section~\ref{sec:c03-conflits}).
\end{enumerate}

\separateur

% ============================================================================
\section{Neutralité et Impartialité}
\label{sec:c03-neutralite}
% ============================================================================

\begin{boxdef}
\termecle{Impartialité}~: disposition de l'enquêteur à suivre les preuves où
elles mènent, sans favoriser une thèse. \termecle{Neutralité}~: absence
d'intérêt personnel dans l'issue. L'expert rapporte ce qu'il établit, y
compris ce qui disculpe.
\end{boxdef}

% ============================================================================
\section{Les Biais}
\label{sec:c03-biais}
% ============================================================================

\begin{tableaucours}{L{3.6cm}L{8.6cm}}{%
    \tableauHeader{Biais} &
    \tableauHeader{Menace}
}
Biais de confirmation & Ne retenir que ce qui confirme la thèse attendue \\
Pression de l'autorité & Orienter l'analyse vers le résultat voulu \\
Attachement à l'outil & Ne voir que ce que l'outil montre \\
Conclusion hâtive & Interpréter avant d'avoir établi \\
\end{tableaucours}
\vspace{-0.8em}
\begin{center}{\small\itshape Tableau~3.1. Biais qui menacent l'enquête}\end{center}

\begin{boxalert}
\textbf{Violation typique.} Chercher dans le disque « la preuve que NDONGO
est coupable ».

\smallskip
\textbf{Règle~:} on cherche ce qui s'est passé, pas une confirmation. On
formule les hypothèses concurrentes et on les éprouve toutes, à charge et à
décharge.
\end{boxalert}

% ============================================================================
\section{Confidentialité et Proportionnalité}
\label{sec:c03-confid}
% ============================================================================

\begin{boxCRO}
\textbf{Confidentialité et proportionnalité.} Un support saisi contient
souvent des données privées sans rapport avec l'affaire (photos, messages
personnels). Les examiner toutes dégrade la \cro{-}{0}{0}~: on limite
l'analyse au périmètre de la réquisition, et l'on protège le reste. La
proportionnalité est une exigence juridique autant qu'éthique.
\end{boxCRO}

% ============================================================================
\section{Conflits d'Intérêts}
\label{sec:c03-conflits}
% ============================================================================

\begin{boxscenario}{Opération TSANAGA. L'expert se déporte-t-il ?}
\textbf{Constatation.} M.~BIKORO constate qu'il connaît personnellement un
cadre de LiZbiZ impliqué dans l'affaire.

\medskip
\textbf{Analyse} (confiance élevée). Un lien personnel peut entacher
l'impartialité, ou l'apparence d'impartialité. M.~BIKORO le signale au
parquet et à Me~TSALA, qui décident de la suite. La transparence sur les
conflits protège l'enquête.
\end{boxscenario}

\begin{boxmetier}{Expert judiciaire}
\textbf{Servir le tribunal, pas une partie.} L'expert judiciaire doit la
vérité au tribunal, non à qui le mandate. Rapporter un élément à décharge
n'est pas une faiblesse~: c'est le cœur de sa mission.
\end{boxmetier}

\separateur

\begin{boxresume}
\textbf{Ce qu'il faut retenir de ce chapitre~:}
\begin{itemize}
    \item \textbf{Impartialité et neutralité}~: suivre les preuves, sans
          intérêt.
    \item \textbf{Les biais} (confirmation, autorité) menacent l'enquête.
    \item \textbf{Proportionnalité}~: limiter l'analyse au périmètre.
    \item \textbf{Conflits d'intérêts}~: les signaler.
    \item \textbf{À charge et à décharge}~: on rapporte tout.
\end{itemize}
\end{boxresume}

\begin{boxinfo}
\textbf{Points de contrôle de l'élève.} Je sais~: définir impartialité et
neutralité~; nommer les biais~; appliquer la proportionnalité~; gérer un
conflit d'intérêts~; rapporter à décharge.
\end{boxinfo}

\bigskip

\begin{boxexo}[title={\ding{45}\quad Exercice~3.1. Impartialité \niveaubase}]
Qu'est-ce que l'impartialité de l'enquêteur~?
\end{boxexo}

\begin{boxexo}[title={\ding{45}\quad Exercice~3.2. Biais \niveaubase}]
Citez deux biais qui menacent une enquête.
\end{boxexo}

\begin{boxexo}[title={\ding{45}\quad Exercice~3.3. À décharge \niveaubase}]
Pourquoi rapporter un élément qui disculpe le mis en cause~?
\end{boxexo}

\begin{boxexo}[title={\ding{45}\quad Exercice~3.4. Confirmation \niveaumoyen}]
En quoi « chercher la preuve de la culpabilité » est-il un biais~?
\end{boxexo}

\begin{boxexo}[title={\ding{45}\quad Exercice~3.5. Proportionnalité \niveaumoyen}]
Un support contient des photos privées sans lien avec l'affaire. Que faire~?
\end{boxexo}

\begin{boxexo}[title={\ding{45}\quad Exercice~3.6. Conflit \niveaumoyen}]
M.~BIKORO connaît une personne impliquée. Comment réagir~?
\end{boxexo}

\begin{boxexo}[title={\ding{45}\quad Exercice~3.7. Pression \niveaumoyen}]
Comment résister à la pression de l'autorité requérante~?
\end{boxexo}

\begin{boxexo}[title={\ding{45}\quad Exercice~3.8. Apparence \niveaumoyen}]
Pourquoi l'apparence d'impartialité compte-t-elle autant que l'impartialité~?
\end{boxexo}

\begin{boxexo}[title={\ding{45}\quad Exercice~3.9. Charte \niveauexpert}]
Rédigez une charte déontologique de l'expert forensique.
\end{boxexo}

\begin{boxexo}[title={\ding{45}\quad Exercice~3.10. Hypothèses \niveauexpert}]
Formulez, pour TSANAGA, deux hypothèses concurrentes et comment les éprouver.
\end{boxexo}

\bigskip

\begin{boxinfo}
\textbf{Pour aller plus loin.}
\begin{itemize}
    \item ISO/IEC 27037~\cite{iso27037}~; ACPO~\cite{acpo2012}.
\end{itemize}
\end{boxinfo}

\begin{boxinfo}
\textbf{Transition.} L'éthique rejoint le droit~: la Partie~II traite du
cadre juridique et de la recevabilité, à commencer par le droit de la preuve
(Chapitre~\ref{chap:droit}).
\end{boxinfo}
}{}

    \part{Le Cadre Juridique et la Recevabilité}
    \label{part:cadre}

        \IfFileExists{partie2/P02_C04_droit_preuve.tex}{% ============================================================================
% Fichier  : partie2/P02_C04_droit_preuve.tex : FORN (LPI342)
% Sources : loi 2010/012, loi 2024/017 ; Malabo 2014 ; Budapest 2001.
% ============================================================================

\chapter{Le Droit de la Preuve Numérique}
\label{chap:droit}
\chaptermark{Le droit de la preuve numérique}

\epigraph{%
    « Une preuve obtenue hors du droit\\
    affaiblit la cause qu'elle prétend servir. »%
}{Principe du cours}

\niveauChapitre{Fondamental}{Chapitre~\ref{chap:ethique}~; notions de droit.}

\begin{boxobjectifs}
\begin{itemize}
    \item Connaître le cadre camerounais de la preuve numérique.
    \item Situer la loi 2010/012 et la loi 2024/017.
    \item Articuler avec Malabo et Budapest.
    \item Comprendre le rôle du juge et de l'expert.
    \item Déduire les règles à respecter pendant l'enquête.
\end{itemize}
\end{boxobjectifs}

\bigskip

% ============================================================================
\section*{Question directrice : que faut-il respecter pour que la preuve tienne devant le juge ?}
% ============================================================================

L'enquête se déroule sous le regard du droit~: le titre pour agir, le respect
des données personnelles, les règles de la preuve. Ce chapitre pose ce cadre.

\medskip
\textbf{Plan du chapitre.}
\begin{enumerate}
    \item Le cadre camerounais (section~\ref{sec:c04-cameroun}).
    \item Les données personnelles (section~\ref{sec:c04-donnees}).
    \item Malabo et Budapest (section~\ref{sec:c04-international}).
    \item Le juge et l'expert (section~\ref{sec:c04-juge}).
\end{enumerate}

\separateur

% ============================================================================
\section{Le Cadre Camerounais}
\label{sec:c04-cameroun}
% ============================================================================

\begin{boxjuris}
La \loicam{loi n°~2010/012 du 21~décembre~2010} relative à la cybersécurité
et à la cybercriminalité réprime l'accès et le maintien frauduleux dans un
système et encadre la preuve électronique. L'accès à un support sans titre
(réquisition, mandat, autorisation) est illicite~; la preuve ainsi obtenue
risque l'irrecevabilité.
\end{boxjuris}

% ============================================================================
\section{Les Données Personnelles}
\label{sec:c04-donnees}
% ============================================================================

\begin{boxjuris}
La \loicam{loi n°~2024/017 du 23~décembre~2024} régit la protection des
données à caractère personnel. Un support d'enquête contient des données
personnelles~: on les traite dans le périmètre de la réquisition, on ne
diffuse rien au-delà, et l'on protège les tiers.
\end{boxjuris}

% ============================================================================
\section{Malabo et Budapest}
\label{sec:c04-international}
% ============================================================================

\begin{tableaucours}{L{3.6cm}L{8.6cm}}{%
    \tableauHeader{Instrument} &
    \tableauHeader{Apport pour la preuve}
}
Convention de Malabo (UA, 2014) & Cadre continental~: infractions, données, coopération \\
Convention de Budapest (2001) & Preuve électronique, perquisition, entraide internationale \\
\end{tableaucours}
\vspace{-0.8em}
\begin{center}{\small\itshape Tableau~4.1. Deux instruments supranationaux~\cite{malabo2014,budapest2001}}\end{center}

% ============================================================================
\section{Le Juge et l'Expert}
\label{sec:c04-juge}
% ============================================================================

Le juge apprécie la preuve~; l'expert l'éclaire. L'expert ne tranche pas la
culpabilité~: il établit des faits techniques et explique leur portée. La
frontière est nette, et la respecter protège la valeur de l'expertise.

\begin{boxscenario}{Opération TSANAGA. Le titre d'abord}
\textbf{Constatation.} Avant toute opération, Me~TSALA vérifie que la
réquisition du parquet couvre bien les supports saisis et le périmètre de
l'analyse.

\medskip
\textbf{Analyse} (confiance élevée). Sans ce contrôle, une analyse hors
périmètre aurait pu être écartée. Le droit n'est pas un obstacle à l'enquête~:
c'est ce qui en garantit la valeur.
\end{boxscenario}

\begin{boxmetier}{Juriste TIC}
\textbf{Sécuriser l'enquête juridiquement.} Le juriste vérifie le titre, le
périmètre, le respect des données, et prépare la recevabilité. Son travail,
en amont, évite que des mois d'enquête ne s'effondrent devant le juge.
\end{boxmetier}

\separateur

\begin{boxresume}
\textbf{Ce qu'il faut retenir de ce chapitre~:}
\begin{itemize}
    \item \textbf{Loi 2010/012}~: accès frauduleux, preuve électronique.
    \item \textbf{Loi 2024/017}~: protection des données personnelles.
    \item \textbf{Malabo et Budapest} encadrent preuve et coopération.
    \item \textbf{Le juge apprécie}, \textbf{l'expert éclaire}~; il ne tranche
          pas la culpabilité.
    \item \textbf{Le titre et le périmètre} se vérifient avant d'agir.
\end{itemize}
\end{boxresume}

\begin{boxinfo}
\textbf{Points de contrôle de l'élève.} Je sais~: citer les textes
camerounais~; situer Malabo et Budapest~; distinguer les rôles du juge et de
l'expert~; énoncer les règles à respecter.
\end{boxinfo}

\bigskip

\begin{boxexo}[title={\ding{45}\quad Exercice~4.1. Le texte \niveaubase}]
Quel texte camerounais encadre la preuve électronique~?
\end{boxexo}

\begin{boxexo}[title={\ding{45}\quad Exercice~4.2. Données \niveaubase}]
Quel texte protège les données personnelles~?
\end{boxexo}

\begin{boxexo}[title={\ding{45}\quad Exercice~4.3. Juge ou expert \niveaubase}]
Qui apprécie la preuve~? Qui l'éclaire~?
\end{boxexo}

\begin{boxexo}[title={\ding{45}\quad Exercice~4.4. Sans titre \niveaumoyen}]
Que risque une preuve obtenue sans titre~?
\end{boxexo}

\begin{boxexo}[title={\ding{45}\quad Exercice~4.5. Périmètre \niveaumoyen}]
Pourquoi vérifier que la réquisition couvre le périmètre d'analyse~?
\end{boxexo}

\begin{boxexo}[title={\ding{45}\quad Exercice~4.6. Budapest \niveaumoyen}]
Qu'apporte la Convention de Budapest à une enquête transfrontalière~?
\end{boxexo}

\begin{boxexo}[title={\ding{45}\quad Exercice~4.7. Tiers \niveaumoyen}]
Comment protéger les tiers dont les données figurent sur un support saisi~?
\end{boxexo}

\begin{boxexo}[title={\ding{45}\quad Exercice~4.8. La frontière \niveaumoyen}]
Pourquoi l'expert ne doit-il pas conclure à la culpabilité~?
\end{boxexo}

\begin{boxexo}[title={\ding{45}\quad Exercice~4.9. Note juridique \niveauexpert}]
Rédigez la note qui sécurise juridiquement l'opération TSANAGA.
\end{boxexo}

\begin{boxexo}[title={\ding{45}\quad Exercice~4.10. Comparer \niveauexpert}]
Comparez le cadre camerounais à la Convention de Malabo sur la preuve.
\end{boxexo}

\bigskip

\begin{boxinfo}
\textbf{Pour aller plus loin.}
\begin{itemize}
    \item Lois 2010/012~\cite{loi2010012} et 2024/017~\cite{loi2024017}~;
          Malabo~\cite{malabo2014}~; Budapest~\cite{budapest2001}.
\end{itemize}
\end{boxinfo}

\begin{boxinfo}
\textbf{Transition.} Le droit exige une continuité de la détention~: le
Chapitre~\ref{chap:custody} traite de la chaîne de possession.
\end{boxinfo}
}{}
        \IfFileExists{partie2/P02_C05_chaine_possession.tex}{% ============================================================================
% Fichier  : partie2/P02_C05_chaine_possession.tex : FORN (LPI342)
% Sources : ISO/IEC 27037 ; ACPO ; RFC 3227. Opération TSANAGA.
% ============================================================================

\chapter{La Chaîne de Possession}
\label{chap:custody}
\chaptermark{La chaîne de possession}

\epigraph{%
    « Une preuve dont on ignore où elle était hier\\
    ne prouve plus rien aujourd'hui. »%
}{Principe du cours}

\niveauChapitre{Fondamental}{Chapitre~\ref{chap:droit}.}

\begin{boxobjectifs}
\begin{itemize}
    \item Définir la chaîne de possession (ou de custody).
    \item Tenir la continuité de la détention d'une pièce.
    \item Sceller, étiqueter, consigner les transferts.
    \item Relier la chaîne de possession à l'opposabilité.
    \item Rédiger une fiche de scellé.
\end{itemize}
\end{boxobjectifs}

\bigskip

% ============================================================================
\section*{Question directrice : comment prouver que le disque analysé est bien celui qui a été saisi ?}
% ============================================================================

Entre la saisie chez M.~NDONGO et l'analyse, le support passe de main en
main. Si l'on ne peut retracer ce parcours, la défense contestera que la
pièce analysée est bien celle saisie. Ce chapitre pose la chaîne de
possession.

\medskip
\textbf{Plan du chapitre.}
\begin{enumerate}
    \item Définir la chaîne (section~\ref{sec:c05-definir}).
    \item Sceller et étiqueter (section~\ref{sec:c05-sceller}).
    \item Consigner les transferts (section~\ref{sec:c05-transferts}).
    \item Chaîne et opposabilité (section~\ref{sec:c05-opposabilite}).
\end{enumerate}

\separateur

% ============================================================================
\section{Définir la Chaîne}
\label{sec:c05-definir}
% ============================================================================

\begin{boxdef}
\termecle{Chaîne de possession} (ou chaîne de custody)~: continuité
documentée de la détention d'une pièce, de la saisie à la présentation,
établissant à tout instant qui la détenait, où, et pour quoi. Toute rupture
(un instant sans détenteur consigné) fragilise la pièce.
\end{boxdef}

% ============================================================================
\section{Sceller et Étiqueter}
\label{sec:c05-sceller}
% ============================================================================

\begin{boxprocedure}
\textbf{Sceller une pièce}
\begin{enumerate}
    \item Identifier la pièce~: numéro de scellé, description, photographie.
    \item Placer sous scellé inviolable, daté et signé.
    \item Calculer l'empreinte (condensat) dès que possible.
    \item Étiqueter~: affaire, pièce, date, agent.
    \item Consigner l'acte dans le registre des scellés.
\end{enumerate}
\end{boxprocedure}

% ============================================================================
\section{Consigner les Transferts}
\label{sec:c05-transferts}
% ============================================================================

\begin{boxobjet}{Scellé n° TSA-2026-07 : clé USB}
\textbf{Nature~:} clé USB saisie au domicile de M.~NDONGO le 20~avril 2026.
\textbf{Scellé~:} TSA-2026-07, inviolable, signé. \textbf{Empreinte~:}
condensat calculé à la réception. \textbf{Transferts~:} parquet $\to$ NYONG
SÉCURITÉ (20/04), coffre (20/04), analyse (22/04). \textbf{Observation~:}
continuité assurée.
\end{boxobjet}

% ============================================================================
\section{Chaîne et Opposabilité}
\label{sec:c05-opposabilite}
% ============================================================================

\begin{boxCRO}
\textbf{La chaîne de possession sert l'opposabilité.} Une pièce à la
détention continue et documentée est opposable \cro{0}{+}{+}~; une pièce dont
la détention présente un trou devient contestable \cro{0}{0}{-}. La chaîne de
possession est la condition première de la recevabilité.
\end{boxCRO}

\begin{boxscenario}{Opération TSANAGA. Un transfert non consigné}
\textbf{Constatation.} La clé USB est restée une nuit sur le bureau d'un
stagiaire, sans consignation.

\medskip
\textbf{Analyse} (confiance élevée). Ce trou, même sans malveillance,
affaiblit l'opposabilité de la clé. M.~BIKORO le consigne honnêtement et
s'appuie davantage sur l'ordinateur, dont la chaîne est intacte. On ne cache
pas une rupture~: on la déclare.
\end{boxscenario}

\begin{boxmetier}{Analyste forensique et expert judiciaire}
\textbf{La chaîne commence à la saisie.} L'expert exige que la chaîne soit
tenue dès la saisie, pas à la réception~: un support mal scellé sur le
terrain arrive déjà fragilisé au laboratoire.
\end{boxmetier}

\separateur

\begin{boxresume}
\textbf{Ce qu'il faut retenir de ce chapitre~:}
\begin{itemize}
    \item \textbf{La chaîne de possession} documente la détention, de la
          saisie à la présentation.
    \item \textbf{On scelle, étiquette, calcule l'empreinte}.
    \item \textbf{On consigne chaque transfert}.
    \item \textbf{Un trou} rend la pièce contestable.
    \item \textbf{On déclare une rupture}, on ne la cache pas.
\end{itemize}
\end{boxresume}

\begin{boxinfo}
\textbf{Points de contrôle de l'élève.} Je sais~: définir la chaîne de
possession~; sceller une pièce~; consigner un transfert~; relier la chaîne à
l'opposabilité~; réagir à une rupture.
\end{boxinfo}

\bigskip

\begin{boxexo}[title={\ding{45}\quad Exercice~5.1. Définir \niveaubase}]
Qu'est-ce que la chaîne de possession~?
\end{boxexo}

\begin{boxexo}[title={\ding{45}\quad Exercice~5.2. Sceller \niveaubase}]
Citez trois éléments d'une fiche de scellé.
\end{boxexo}

\begin{boxexo}[title={\ding{45}\quad Exercice~5.3. L'empreinte \niveaubase}]
Pourquoi calculer l'empreinte dès la réception~?
\end{boxexo}

\begin{boxexo}[title={\ding{45}\quad Exercice~5.4. Le trou \niveaumoyen}]
La clé USB est restée une nuit sans consignation. Quel effet~?
\end{boxexo}

\begin{boxexo}[title={\ding{45}\quad Exercice~5.5. Déclarer \niveaumoyen}]
Pourquoi déclarer une rupture plutôt que la cacher~?
\end{boxexo}

\begin{boxexo}[title={\ding{45}\quad Exercice~5.6. CRO \niveaumoyen}]
Expliquez l'analyse CRO d'une chaîne continue et d'une chaîne rompue.
\end{boxexo}

\begin{boxexo}[title={\ding{45}\quad Exercice~5.7. Terrain \niveaumoyen}]
Pourquoi la chaîne doit-elle commencer à la saisie, pas à la réception~?
\end{boxexo}

\begin{boxexo}[title={\ding{45}\quad Exercice~5.8. Identité de la pièce \niveaumoyen}]
Comment prouver que le disque analysé est bien celui saisi~?
\end{boxexo}

\begin{boxexo}[title={\ding{45}\quad Exercice~5.9. Fiche de scellé \niveauexpert}]
Rédigez la fiche de scellé complète de la clé USB TSA-2026-07.
\end{boxexo}

\begin{boxexo}[title={\ding{45}\quad Exercice~5.10. Registre \niveauexpert}]
Concevez le registre des scellés du cabinet NYONG SÉCURITÉ.
\end{boxexo}

\bigskip

\begin{boxinfo}
\textbf{Pour aller plus loin.}
\begin{itemize}
    \item ISO/IEC 27037~\cite{iso27037}~; ACPO~\cite{acpo2012}~;
          RFC~3227~\cite{rfc3227}.
\end{itemize}
\end{boxinfo}

\begin{boxinfo}
\textbf{Transition.} Chaîne et droit convergent sur la question de la
recevabilité~: le Chapitre~\ref{chap:recevabilite} la traite.
\end{boxinfo}
}{}
        \IfFileExists{partie2/P02_C06_recevabilite.tex}{% ============================================================================
% Fichier  : partie2/P02_C06_recevabilite.tex : FORN (LPI342)
% Sources : ISO/IEC 27037, 27042 ; ACPO ; principes de la preuve.
% ============================================================================

\chapter{Recevabilité et Admissibilité}
\label{chap:recevabilite}
\chaptermark{Recevabilité et admissibilité}

\epigraph{%
    « Une vérité que l'on ne peut pas défendre\\
    n'entre pas dans un prétoire. »%
}{Principe du cours}

\niveauChapitre{Intermédiaire}{Chapitres~\ref{chap:droit} et~\ref{chap:custody}.}

\begin{boxobjectifs}
\begin{itemize}
    \item Définir la recevabilité d'une preuve numérique.
    \item Connaître ses conditions (licéité, intégrité, authenticité, traçabilité).
    \item Comprendre les quatre principes de l'ACPO.
    \item Relier la recevabilité au Trilemme CRO.
    \item Reconnaître ce qui rend une preuve irrecevable.
\end{itemize}
\end{boxobjectifs}

\bigskip

% ============================================================================
\section*{Question directrice : qu'est-ce qui rend une preuve numérique recevable, ou irrecevable, devant la justice ?}
% ============================================================================

Une preuve peut être exacte et pourtant écartée~: obtenue sans titre,
altérée, non traçable. Ce chapitre pose les conditions de la recevabilité et
les principes qui la garantissent.

\medskip
\textbf{Plan du chapitre.}
\begin{enumerate}
    \item Définir la recevabilité (section~\ref{sec:c06-definir}).
    \item Les conditions (section~\ref{sec:c06-conditions}).
    \item Les principes de l'ACPO (section~\ref{sec:c06-acpo}).
    \item Ce qui rend irrecevable (section~\ref{sec:c06-irrecevable}).
\end{enumerate}

\separateur

% ============================================================================
\section{Définir la Recevabilité}
\label{sec:c06-definir}
% ============================================================================

\begin{boxdef}
\termecle{Recevabilité}~: aptitude d'une preuve à être admise et prise en
compte par la justice. Elle tient à la \emph{manière} dont la preuve a été
obtenue, préservée et documentée, non à sa seule exactitude.
\end{boxdef}

% ============================================================================
\section{Les Conditions}
\label{sec:c06-conditions}
% ============================================================================

\begin{tableaucours}{L{3.2cm}L{9.0cm}}{%
    \tableauHeader{Condition} &
    \tableauHeader{Exigence}
}
Licéité & Obtenue avec un titre (réquisition, mandat, autorisation) \\
Intégrité & Non altérée depuis l'acquisition (condensat) \\
Authenticité & Bien ce qu'elle prétend être, de la source déclarée \\
Traçabilité & Chaîne de possession continue, actes documentés \\
\end{tableaucours}
\vspace{-0.8em}
\begin{center}{\small\itshape Tableau~6.1. Conditions de la recevabilité}\end{center}

% ============================================================================
\section{Les Principes de l'ACPO}
\label{sec:c06-acpo}
% ============================================================================

\begin{boxprocedure}
\textbf{Les quatre principes de l'ACPO}
\begin{enumerate}
    \item Ne pas modifier les données originales.
    \item En cas d'accès nécessaire à l'original, être compétent et expliquer
          l'incidence de ses actes.
    \item Tracer tous les actes~: un tiers doit pouvoir les reproduire et
          obtenir le même résultat.
    \item Le responsable de l'enquête garantit le respect du droit et de ces
          principes.
\end{enumerate}
\end{boxprocedure}

% ============================================================================
\section{Ce qui Rend Irrecevable}
\label{sec:c06-irrecevable}
% ============================================================================

\begin{boxCRO}
\textbf{La recevabilité est l'opposabilité.} La dimension O du Trilemme est la
recevabilité~; la F (fiabilité) recouvre intégrité et authenticité~; la C
(confidentialité) impose la proportionnalité. Une preuve bien conduite vise
\cro{0}{+}{+}.
\end{boxCRO}

\begin{boxscenario}{Opération TSANAGA. Deux pièces, deux sorts}
\textbf{Constatation.} L'ordinateur, acquis selon l'ISO/IEC 27037, à la
chaîne continue, est recevable. La clé USB, dont la chaîne présente un trou
(Chapitre~\ref{chap:custody}), est fragilisée.

\medskip
\textbf{Analyse} (confiance élevée). Les faits tirés de l'ordinateur tiendront~;
ceux tirés de la clé seule seront discutés. On construit la démonstration sur
les pièces les plus solides, et l'on recoupe.
\end{boxscenario}

\begin{boxmetier}{Expert judiciaire et avocat}
\textbf{Anticiper la contestation.} L'expert et l'avocat se demandent, pour
chaque pièce~: « comment la défense pourrait-elle la contester~? ». Ils
corrigent en amont ce qui peut l'être, et signalent honnêtement le reste.
\end{boxmetier}

\separateur

\begin{boxresume}
\textbf{Ce qu'il faut retenir de ce chapitre~:}
\begin{itemize}
    \item \textbf{La recevabilité} tient à la manière, pas à la seule
          exactitude.
    \item \textbf{Quatre conditions}~: licéité, intégrité, authenticité,
          traçabilité.
    \item \textbf{Quatre principes de l'ACPO}~: ne pas altérer, être compétent,
          tracer, garantir.
    \item \textbf{L'opposabilité} du CRO \emph{est} la recevabilité.
    \item \textbf{On anticipe} la contestation.
\end{itemize}
\end{boxresume}

\begin{boxinfo}
\textbf{Points de contrôle de l'élève.} Je sais~: définir la recevabilité~;
citer ses conditions~; énoncer les principes de l'ACPO~; relier au CRO~;
reconnaître une preuve irrecevable.
\end{boxinfo}

\bigskip

\begin{boxexo}[title={\ding{45}\quad Exercice~6.1. Définir \niveaubase}]
À quoi tient la recevabilité d'une preuve~?
\end{boxexo}

\begin{boxexo}[title={\ding{45}\quad Exercice~6.2. Conditions \niveaubase}]
Citez les quatre conditions de la recevabilité.
\end{boxexo}

\begin{boxexo}[title={\ding{45}\quad Exercice~6.3. ACPO \niveaubase}]
Énoncez le premier principe de l'ACPO.
\end{boxexo}

\begin{boxexo}[title={\ding{45}\quad Exercice~6.4. Exacte mais écartée \niveaumoyen}]
Comment une preuve exacte peut-elle être écartée~?
\end{boxexo}

\begin{boxexo}[title={\ding{45}\quad Exercice~6.5. Reproductible \niveaumoyen}]
Pourquoi un tiers doit-il pouvoir reproduire l'analyse~?
\end{boxexo}

\begin{boxexo}[title={\ding{45}\quad Exercice~6.6. Deux pièces \niveaumoyen}]
Pourquoi l'ordinateur est-il recevable et la clé USB fragilisée~?
\end{boxexo}

\begin{boxexo}[title={\ding{45}\quad Exercice~6.7. CRO \niveaumoyen}]
Pourquoi dit-on que l'opposabilité est la recevabilité~?
\end{boxexo}

\begin{boxexo}[title={\ding{45}\quad Exercice~6.8. Anticiper \niveaumoyen}]
Comment anticiper la contestation d'une pièce par la défense~?
\end{boxexo}

\begin{boxexo}[title={\ding{45}\quad Exercice~6.9. Liste de contrôle \niveauexpert}]
Rédigez une liste de contrôle de recevabilité pour TSANAGA.
\end{boxexo}

\begin{boxexo}[title={\ding{45}\quad Exercice~6.10. Note de recevabilité \niveauexpert}]
Rédigez la note qui établit la recevabilité de l'ordinateur saisi.
\end{boxexo}

\bigskip

\begin{boxinfo}
\textbf{Pour aller plus loin.}
\begin{itemize}
    \item ISO/IEC 27037~\cite{iso27037}, 27042~\cite{iso27042}~;
          ACPO~\cite{acpo2012}.
\end{itemize}
\end{boxinfo}

\begin{boxinfo}
\textbf{Transition.} La recevabilité se joue dès l'acquisition~: la
Partie~III s'ouvre sur l'identification et la préservation des sources
(Chapitre~\ref{chap:preserver}).
\end{boxinfo}
}{}

    \part{L'Acquisition des Preuves}
    \label{part:acquisition}

        \IfFileExists{partie3/P03_C07_identifier_preserver.tex}{% ============================================================================
% Fichier  : partie3/P03_C07_identifier_preserver.tex : FORN (LPI342)
% Sources : ISO/IEC 27037 ; RFC 3227 ; ACPO. Opération TSANAGA.
% ============================================================================

\chapter{Identifier et Préserver les Sources}
\label{chap:preserver}
\chaptermark{Identifier et préserver}

\epigraph{%
    « La preuve la plus solide est celle que l'on n'a pas abîmée\\
    en voulant la regarder. »%
}{Principe du cours}

\niveauChapitre{Fondamental}{Partie~II.}

\begin{boxobjectifs}
\begin{itemize}
    \item Identifier les sources potentielles de preuve.
    \item Préserver une scène numérique sans l'altérer.
    \item Décider d'éteindre, de débrancher ou de capturer à chaud.
    \item Protéger contre l'écriture (bloqueur).
    \item Documenter l'état initial.
\end{itemize}
\end{boxobjectifs}

\bigskip

% ============================================================================
\section*{Question directrice : devant l'ordinateur de M. NDONGO, faut-il l'éteindre, le débrancher, ou le laisser allumé ?}
% ============================================================================

Chaque geste sur un système vivant a un coût. Ce chapitre apprend à
identifier les sources, à préserver la scène et à décider de la conduite à
tenir selon l'état du support.

\medskip
\textbf{Plan du chapitre.}
\begin{enumerate}
    \item Les sources de preuve (section~\ref{sec:c07-sources}).
    \item Préserver la scène (section~\ref{sec:c07-scene}).
    \item Éteindre, débrancher ou capturer (section~\ref{sec:c07-decision}).
    \item Protéger contre l'écriture (section~\ref{sec:c07-bloqueur}).
\end{enumerate}

\separateur

% ============================================================================
\section{Les Sources de Preuve}
\label{sec:c07-sources}
% ============================================================================

\begin{tableaucours}{L{3.6cm}L{8.6cm}}{%
    \tableauHeader{Source} &
    \tableauHeader{Ce qu'elle peut contenir}
}
Disques et clés USB & Fichiers, fragments, artefacts, fichiers supprimés \\
Mémoire vive & Processus, secrets, données non écrites \\
Téléphones & Messages, applications, géolocalisation \\
Équipements réseau & Journaux, flux, configurations \\
Services en ligne & Comptes, journaux, sauvegardes \\
\end{tableaucours}
\vspace{-0.8em}
\begin{center}{\small\itshape Tableau~7.1. Sources potentielles de preuve}\end{center}

% ============================================================================
\section{Préserver la Scène}
\label{sec:c07-scene}
% ============================================================================

\begin{boxprocedure}
\textbf{Préserver une scène numérique}
\begin{enumerate}
    \item Photographier l'état initial (écran, branchements, voyants).
    \item Ne toucher à rien avant d'avoir décidé et documenté.
    \item Isoler des réseaux (éviter effacement ou verrouillage à distance).
    \item Noter l'heure du système et l'heure réelle (écart).
    \item Consigner chaque geste.
\end{enumerate}
\end{boxprocedure}

% ============================================================================
\section{Éteindre, Débrancher ou Capturer}
\label{sec:c07-decision}
% ============================================================================

\begin{tableaucours}{L{3.2cm}L{4.6cm}L{4.4cm}}{%
    \tableauHeader{État} &
    \tableauHeader{Conduite} &
    \tableauHeader{Raison}
}
Éteint & Ne pas allumer, acquérir le disque & L'allumage modifie le disque \\
Allumé, données volatiles utiles & Capturer la mémoire d'abord & La mémoire meurt à l'extinction \\
Allumé, risque de chiffrement & Éviter l'extinction brutale & Peut déclencher un verrouillage \\
\end{tableaucours}
\vspace{-0.8em}
\begin{center}{\small\itshape Tableau~7.2. Décider selon l'état du support}\end{center}

\begin{boxalert}
\textbf{Violation typique.} Allumer l'ordinateur éteint « juste pour voir ce
qu'il contient ».

\smallskip
\textbf{Règle~:} allumer modifie des dizaines d'artefacts (dates d'accès,
démarrage). On n'allume pas un support éteint~: on acquiert son disque. On ne
touche au vivant que pour une raison documentée.
\end{boxalert}

% ============================================================================
\section{Protéger contre l'Écriture}
\label{sec:c07-bloqueur}
% ============================================================================

\begin{boxscenario}{Opération TSANAGA. Le bloqueur en écriture}
\textbf{Constatation.} Pour lire le disque et la clé saisis, Mme~DJUIDJE passe
par un \termecle{bloqueur en écriture} (matériel ou logiciel), qui autorise la
lecture mais interdit toute modification du support.

\medskip
\textbf{Analyse} (confiance élevée). Sans bloqueur, le simple branchement
aurait pu écrire sur le support (horodatage, index). Le bloqueur garantit
l'intégrité exigée pour la recevabilité (Chapitre~\ref{chap:recevabilite}).
\end{boxscenario}

\begin{boxmetier}{Analyste forensique}
\textbf{Lire sans écrire.} L'analyste branche toujours un support de preuve
derrière un bloqueur en écriture, et vérifie par condensat qu'il n'a rien
modifié. C'est un réflexe, pas une option.
\end{boxmetier}

\separateur

\begin{boxresume}
\textbf{Ce qu'il faut retenir de ce chapitre~:}
\begin{itemize}
    \item \textbf{Les sources}~: disques, mémoire, téléphones, réseau,
          services.
    \item \textbf{Préserver la scène}~: photographier, isoler, ne rien
          toucher sans raison.
    \item \textbf{Éteint}~: ne pas allumer. \textbf{Allumé}~: capturer la
          mémoire d'abord.
    \item \textbf{Un bloqueur en écriture} protège l'intégrité.
    \item \textbf{On documente} l'état initial et chaque geste.
\end{itemize}
\end{boxresume}

\begin{boxinfo}
\textbf{Points de contrôle de l'élève.} Je sais~: citer les sources~;
préserver une scène~; décider selon l'état~; utiliser un bloqueur~;
documenter l'état initial.
\end{boxinfo}

\bigskip

\begin{boxexo}[title={\ding{45}\quad Exercice~7.1. Sources \niveaubase}]
Citez quatre sources potentielles de preuve.
\end{boxexo}

\begin{boxexo}[title={\ding{45}\quad Exercice~7.2. Éteint \niveaubase}]
Que faire d'un ordinateur éteint~? Pourquoi~?
\end{boxexo}

\begin{boxexo}[title={\ding{45}\quad Exercice~7.3. Bloqueur \niveaubase}]
À quoi sert un bloqueur en écriture~?
\end{boxexo}

\begin{boxexo}[title={\ding{45}\quad Exercice~7.4. Allumé \niveaumoyen}]
Un ordinateur est allumé avec des données volatiles utiles. Que faire~?
\end{boxexo}

\begin{boxexo}[title={\ding{45}\quad Exercice~7.5. Allumer \niveaumoyen}]
Pourquoi n'allume-t-on pas un support éteint pour « voir »~?
\end{boxexo}

\begin{boxexo}[title={\ding{45}\quad Exercice~7.6. Isoler \niveaumoyen}]
Pourquoi isoler un support des réseaux dès la saisie~?
\end{boxexo}

\begin{boxexo}[title={\ding{45}\quad Exercice~7.7. Chiffrement \niveaumoyen}]
Pourquoi éviter l'extinction brutale d'un système peut-être chiffré~?
\end{boxexo}

\begin{boxexo}[title={\ding{45}\quad Exercice~7.8. État initial \niveaumoyen}]
Que consigne-t-on de l'état initial d'une scène~?
\end{boxexo}

\begin{boxexo}[title={\ding{45}\quad Exercice~7.9. Procédure de saisie \niveauexpert}]
Rédigez la procédure de saisie et de préservation pour TSANAGA.
\end{boxexo}

\begin{boxexo}[title={\ding{45}\quad Exercice~7.10. Fiche de décision \niveauexpert}]
Concevez une fiche d'aide à la décision « éteint / allumé / chiffré ».
\end{boxexo}

\bigskip

\begin{boxinfo}
\textbf{Pour aller plus loin.}
\begin{itemize}
    \item ISO/IEC 27037~\cite{iso27037}~; RFC~3227~\cite{rfc3227}~;
          ACPO~\cite{acpo2012}.
\end{itemize}
\end{boxinfo}

\begin{boxinfo}
\textbf{Transition.} Une fois le support préservé, on le duplique~: le
Chapitre~\ref{chap:duplication} traite de la duplication et du hachage.
\end{boxinfo}
}{}
        \IfFileExists{partie3/P03_C08_duplication_hachage.tex}{% ============================================================================
% Fichier  : partie3/P03_C08_duplication_hachage.tex : FORN (LPI342)
% Sources : ISO/IEC 27037 ; NIST SP 800-86 ; ACPO. Opération TSANAGA.
% ============================================================================

\chapter{Duplication et Hachage}
\label{chap:duplication}
\chaptermark{Duplication et hachage}

\epigraph{%
    « Le condensat est la signature de l'intégrité~:\\
    un bit change, il change tout. »%
}{Principe du cours}

\niveauChapitre{Intermédiaire}{Chapitre~\ref{chap:preserver}.}

\begin{boxobjectifs}
\begin{itemize}
    \item Comprendre la copie bit à bit (image).
    \item Comprendre les fonctions de hachage (MD5, SHA-256).
    \item Utiliser le condensat pour prouver l'intégrité.
    \item Travailler sur la copie, jamais sur l'original.
    \item Documenter l'acquisition.
\end{itemize}
\end{boxobjectifs}

\bigskip

% ============================================================================
\section*{Question directrice : comment prouver que la copie du disque est identique à l'original, et le restera ?}
% ============================================================================

On n'analyse jamais l'original~: on en fait une image, dont on prouve la
fidélité par un condensat. Ce chapitre explique la duplication et le hachage,
cœur technique de l'intégrité.

\medskip
\textbf{Plan du chapitre.}
\begin{enumerate}
    \item La copie bit à bit (section~\ref{sec:c08-image}).
    \item Les fonctions de hachage (section~\ref{sec:c08-hachage}).
    \item Prouver l'intégrité (section~\ref{sec:c08-integrite}).
    \item Documenter l'acquisition (section~\ref{sec:c08-documenter}).
\end{enumerate}

\separateur

% ============================================================================
\section{La Copie Bit à Bit}
\label{sec:c08-image}
% ============================================================================

\begin{boxdef}
\termecle{Image forensique}~: copie \emph{bit à bit} d'un support, qui
reproduit non seulement les fichiers visibles, mais aussi l'espace non alloué,
les fragments et les métadonnées. Elle se distingue d'une simple copie de
fichiers, qui perd l'essentiel des traces.
\end{boxdef}

% ============================================================================
\section{Les Fonctions de Hachage}
\label{sec:c08-hachage}
% ============================================================================

\begin{boxdef}
\termecle{Fonction de hachage}~: fonction qui transforme une donnée, quelle
que soit sa taille, en une empreinte de longueur fixe (le \termecle{condensat}).
Un changement d'un seul bit produit un condensat entièrement différent. MD5 et
SHA-256 en sont des exemples~; SHA-256 est préféré, MD5 étant considéré comme
faible pour les usages sensibles.
\end{boxdef}

\begin{tableaucours}{L{3.0cm}L{9.2cm}}{%
    \tableauHeader{Fonction} &
    \tableauHeader{Usage}
}
MD5 & Historique, rapide~; éviter seule pour un usage sensible \\
SHA-256 & Recommandée pour l'intégrité forensique \\
\end{tableaucours}
\vspace{-0.8em}
\begin{center}{\small\itshape Tableau~8.1. Deux fonctions de hachage}\end{center}

% ============================================================================
\section{Prouver l'Intégrité}
\label{sec:c08-integrite}
% ============================================================================

On calcule le condensat de l'original \emph{avant} la copie, puis celui de la
copie. S'ils sont identiques, la copie est fidèle et l'original n'a pas
changé. On recalcule le condensat de la copie avant chaque usage~: s'il
correspond toujours, la copie est restée intègre.

\begin{boxCRO}
\textbf{Le condensat sert l'opposabilité.} Un condensat calculé et consigné
établit la \cro{0}{+}{+}~: il prouve l'intégrité (fiabilité) et rend la pièce
défendable (opposabilité). Sans condensat, l'intégrité n'est qu'affirmée.
\end{boxCRO}

% ============================================================================
\section{Documenter l'Acquisition}
\label{sec:c08-documenter}
% ============================================================================

\begin{boxscenario}{Opération TSANAGA. L'image du disque}
\textbf{Constatation.} Derrière un bloqueur en écriture, Mme~NKENG réalise
une image bit à bit du disque de l'ordinateur et calcule son condensat
SHA-256, qu'elle consigne. Toute l'analyse se fera sur cette image.

\medskip
\textbf{Analyse} (confiance élevée). Le condensat consigné permettra, à tout
moment, de prouver que l'image analysée est fidèle à l'original. L'original
est remis sous scellé, intact.
\end{boxscenario}

\begin{boxobjet}{Image disque TSA-2026-01}
\textbf{Nature~:} image bit à bit du disque de l'ordinateur saisi.
\textbf{Source~:} disque TSA-2026-01, derrière bloqueur en écriture.
\textbf{Intégrité~:} condensat SHA-256 consigné à l'acquisition.
\textbf{Opposabilité~:} chaîne de possession continue. \textbf{Usage~:}
analyse sur l'image, original scellé.
\end{boxobjet}

\begin{boxmetier}{Analyste forensique}
\textbf{Hacher, copier, re-hacher.} L'analyste calcule le condensat avant et
après la copie, et le recalcule avant chaque session de travail. Trois
gestes, une garantie d'intégrité à toute épreuve.
\end{boxmetier}

\separateur

\begin{boxresume}
\textbf{Ce qu'il faut retenir de ce chapitre~:}
\begin{itemize}
    \item \textbf{L'image bit à bit} copie tout, y compris l'espace non
          alloué.
    \item \textbf{Le hachage} produit un condensat sensible au moindre bit.
    \item \textbf{SHA-256} est préféré à MD5 pour l'intégrité.
    \item \textbf{Le condensat} prouve la fidélité de la copie.
    \item \textbf{On analyse la copie}, l'original reste scellé.
\end{itemize}
\end{boxresume}

\begin{boxinfo}
\textbf{Points de contrôle de l'élève.} Je sais~: définir l'image bit à bit~;
expliquer le hachage~; prouver l'intégrité par condensat~; documenter une
acquisition.
\end{boxinfo}

\bigskip

\begin{boxexo}[title={\ding{45}\quad Exercice~8.1. Image \niveaubase}]
Qu'est-ce qu'une image bit à bit, et en quoi diffère-t-elle d'une copie de
fichiers~?
\end{boxexo}

\begin{boxexo}[title={\ding{45}\quad Exercice~8.2. Condensat \niveaubase}]
Que se passe-t-il pour le condensat si un seul bit change~?
\end{boxexo}

\begin{boxexo}[title={\ding{45}\quad Exercice~8.3. MD5 ou SHA-256 \niveaubase}]
Quelle fonction préférer pour l'intégrité forensique~? Pourquoi~?
\end{boxexo}

\begin{boxexo}[title={\ding{45}\quad Exercice~8.4. Prouver \niveaumoyen}]
Comment prouver que la copie est fidèle et le reste~?
\end{boxexo}

\begin{boxexo}[title={\ding{45}\quad Exercice~8.5. Avant la copie \niveaumoyen}]
Pourquoi calculer le condensat de l'original avant la copie~?
\end{boxexo}

\begin{boxexo}[title={\ding{45}\quad Exercice~8.6. L'espace non alloué \niveaumoyen}]
Pourquoi une image bit à bit capture-t-elle l'espace non alloué~?
\end{boxexo}

\begin{boxexo}[title={\ding{45}\quad Exercice~8.7. CRO \niveaumoyen}]
Expliquez l'analyse CRO du condensat.
\end{boxexo}

\begin{boxexo}[title={\ding{45}\quad Exercice~8.8. Re-hacher \niveaumoyen}]
Pourquoi recalculer le condensat avant chaque session de travail~?
\end{boxexo}

\begin{boxexo}[title={\ding{45}\quad Exercice~8.9. Procédure \niveauexpert}]
Rédigez la procédure d'acquisition par image du disque de TSANAGA.
\end{boxexo}

\begin{boxexo}[title={\ding{45}\quad Exercice~8.10. Fiche d'image \niveauexpert}]
Concevez la fiche d'acquisition (image) du cabinet NYONG SÉCURITÉ.
\end{boxexo}

\bigskip

\begin{boxinfo}
\textbf{Pour aller plus loin.}
\begin{itemize}
    \item ISO/IEC 27037~\cite{iso27037}~; NIST SP~800-86~\cite{nist80086}.
\end{itemize}
\end{boxinfo}

\begin{boxinfo}
\textbf{Transition.} Toutes les traces ne survivent pas également~: le
Chapitre~\ref{chap:acquisition} traite de l'ordre de volatilité et de
l'acquisition à chaud.
\end{boxinfo}
}{}
        \IfFileExists{partie3/P03_C09_volatilite_acquisition.tex}{% ============================================================================
% Fichier  : partie3/P03_C09_volatilite_acquisition.tex : FORN (LPI342)
% Sources : RFC 3227 ; NIST SP 800-86 ; ISO/IEC 27037. Opération TSANAGA.
% ============================================================================

\chapter{Ordre de Volatilité et Acquisition à Chaud}
\label{chap:acquisition}
\chaptermark{Volatilité et acquisition à chaud}

\epigraph{%
    « On capture d'abord ce qui s'efface le plus vite. »%
}{Principe du cours}

\niveauChapitre{Intermédiaire à expert}{Chapitre~\ref{chap:duplication}.}

\begin{boxobjectifs}
\begin{itemize}
    \item Énoncer l'ordre de volatilité.
    \item Décider quand acquérir à chaud.
    \item Capturer la mémoire vive et l'état du système.
    \item Mesurer l'empreinte de l'acquisition à chaud.
    \item Documenter un arbitrage CRO.
\end{itemize}
\end{boxobjectifs}

\bigskip

% ============================================================================
\section*{Question directrice : sur un système allumé, par quoi commencer l'acquisition ?}
% ============================================================================

Sur un système vivant, certaines traces (mémoire, connexions) disparaissent à
l'extinction, d'autres (disque) sont stables. L'ordre de volatilité fixe la
priorité. Ce chapitre l'applique.

\medskip
\textbf{Plan du chapitre.}
\begin{enumerate}
    \item L'ordre de volatilité (section~\ref{sec:c09-volatilite}).
    \item L'acquisition à chaud (section~\ref{sec:c09-chaud}).
    \item Capturer la mémoire (section~\ref{sec:c09-memoire}).
    \item L'arbitrage CRO (section~\ref{sec:c09-cro}).
\end{enumerate}

\separateur

% ============================================================================
\section{L'Ordre de Volatilité}
\label{sec:c09-volatilite}
% ============================================================================

\begin{boxprocedure}
\textbf{Acquérir selon l'ordre de volatilité (RFC 3227)}
\begin{enumerate}
    \item Mémoire vive et état (processus, connexions, utilisateurs).
    \item Fichiers temporaires, données de session.
    \item Disque (image bit à bit).
    \item Journaux distants, configurations.
    \item Supports d'archives.
\end{enumerate}
\end{boxprocedure}

% ============================================================================
\section{L'Acquisition à Chaud}
\label{sec:c09-chaud}
% ============================================================================

\begin{boxdef}
\termecle{Acquisition à chaud}~: capture de données sur un système en
fonctionnement, avant son extinction. Nécessaire pour la mémoire et l'état
vivant, elle laisse inévitablement une empreinte sur le système, qu'on
minimise et documente.
\end{boxdef}

\begin{boxalert}
\textbf{Violation typique.} Éteindre un système allumé « pour le mettre sous
scellé » sans capturer la mémoire.

\smallskip
\textbf{Règle~:} si des données volatiles sont utiles (session chiffrée
ouverte, connexions, secrets en clair), on les capture d'abord. Éteindre sans
réfléchir détruit une preuve unique.
\end{boxalert}

% ============================================================================
\section{Capturer la Mémoire}
\label{sec:c09-memoire}
% ============================================================================

On capture la mémoire avec un outil à empreinte minimale, documenté, en
écrivant sur un support externe, puis on calcule le condensat de la capture.
La mémoire révèle ce que le disque ignore~: processus, connexions, clés de
chiffrement en clair.

% ============================================================================
\section{L'Arbitrage CRO}
\label{sec:c09-cro}
% ============================================================================

\begin{boxCRO}
\textbf{L'acquisition à chaud, un arbitrage.} Toucher au vivant laisse une
empreinte~: \cro{0}{-}{+}. On accepte une légère altération (fiabilité) pour
sauver une trace unique et opposable (opposabilité), et on la documente pour
la rendre défendable. Le choix se justifie et se trace.
\end{boxCRO}

\begin{boxscenario}{Opération TSANAGA. Une session chiffrée ouverte}
\textbf{Constatation.} À la saisie, l'ordinateur de M.~NDONGO était allumé,
avec un conteneur chiffré monté (déchiffré en mémoire).

\medskip
\textbf{Analyse} (confiance élevée). Éteindre aurait reverrouillé le
conteneur, rendant l'accès très difficile. L'équipe a capturé la mémoire à
chaud (où la clé était en clair), puis acquis le disque. L'arbitrage, risqué,
a été documenté et justifié.
\end{boxscenario}

\begin{boxmetier}{Analyste forensique}
\textbf{Décider vite, tracer tout.} Sur une scène vivante, l'analyste décide
en minutes, mais documente chaque choix~: pourquoi capturer à chaud, avec
quel outil, à quelle heure. La trace de la décision vaut autant que la
capture.
\end{boxmetier}

\separateur

\begin{boxresume}
\textbf{Ce qu'il faut retenir de ce chapitre~:}
\begin{itemize}
    \item \textbf{L'ordre de volatilité}~: mémoire, puis disque, puis
          archives.
    \item \textbf{L'acquisition à chaud} sauve les traces volatiles, en
          laissant une empreinte.
    \item \textbf{On capture la mémoire} avant l'extinction, si elle est
          utile.
    \item \textbf{L'arbitrage CRO} se justifie et se documente.
    \item \textbf{Une session chiffrée ouverte} change la décision.
\end{itemize}
\end{boxresume}

\begin{boxinfo}
\textbf{Points de contrôle de l'élève.} Je sais~: énoncer l'ordre de
volatilité~; décider d'acquérir à chaud~; capturer la mémoire~; analyser le
CRO de l'acquisition~; documenter l'arbitrage.
\end{boxinfo}

\bigskip

\begin{boxexo}[title={\ding{45}\quad Exercice~9.1. Volatilité \niveaubase}]
Classez du plus volatil au plus stable~: disque, mémoire, archives.
\end{boxexo}

\begin{boxexo}[title={\ding{45}\quad Exercice~9.2. À chaud \niveaubase}]
Qu'est-ce qu'une acquisition à chaud~?
\end{boxexo}

\begin{boxexo}[title={\ding{45}\quad Exercice~9.3. Support \niveaubase}]
Sur quel support écrit-on une capture mémoire~?
\end{boxexo}

\begin{boxexo}[title={\ding{45}\quad Exercice~9.4. Éteindre \niveaumoyen}]
Pourquoi ne pas éteindre un système allumé avant de capturer la mémoire~?
\end{boxexo}

\begin{boxexo}[title={\ding{45}\quad Exercice~9.5. Conteneur chiffré \niveaumoyen}]
Pourquoi la session chiffrée ouverte change-t-elle la décision~?
\end{boxexo}

\begin{boxexo}[title={\ding{45}\quad Exercice~9.6. CRO \niveaumoyen}]
Expliquez l'arbitrage CRO de l'acquisition à chaud.
\end{boxexo}

\begin{boxexo}[title={\ding{45}\quad Exercice~9.7. Empreinte \niveaumoyen}]
Pourquoi documenter l'empreinte laissée par l'outil d'acquisition~?
\end{boxexo}

\begin{boxexo}[title={\ding{45}\quad Exercice~9.8. Ce que révèle la mémoire \niveaumoyen}]
Citez trois éléments que la mémoire révèle et que le disque ignore.
\end{boxexo}

\begin{boxexo}[title={\ding{45}\quad Exercice~9.9. Procédure \niveauexpert}]
Rédigez la procédure d'acquisition à chaud de l'ordinateur de TSANAGA.
\end{boxexo}

\begin{boxexo}[title={\ding{45}\quad Exercice~9.10. Justification \niveauexpert}]
Rédigez la note qui justifie l'acquisition à chaud devant le tribunal.
\end{boxexo}

\bigskip

\begin{boxinfo}
\textbf{Pour aller plus loin.}
\begin{itemize}
    \item RFC~3227~\cite{rfc3227}~; NIST SP~800-86~\cite{nist80086}~;
          ISO/IEC 27037~\cite{iso27037}.
\end{itemize}
\end{boxinfo}

\begin{boxinfo}
\textbf{Transition.} L'acquisition faite, l'analyse commence~: la Partie~IV
parcourt les domaines, à commencer par le disque
(Chapitre~\ref{chap:disque}).
\end{boxinfo}
}{}

    \part{L'Analyse par Domaine}
    \label{part:analyse}

        \IfFileExists{partie4/P04_C10_disque_systemes_fichiers.tex}{% ============================================================================
% Fichier  : partie4/P04_C10_disque_systemes_fichiers.tex : FORN (LPI342)
% Sources : ISO/IEC 27042 ; NIST SP 800-86 ; documentation FAT, NTFS, EXT.
% ============================================================================

\chapter{Forensique Disque et Systèmes de Fichiers}
\label{chap:disque}
\chaptermark{Forensique disque}

\epigraph{%
    « Le système de fichiers est une carte~;\\
    l'enquêteur apprend à lire ses marges autant que son centre. »%
}{Principe du cours}

\niveauChapitre{Intermédiaire à expert}{Partie~III~; administration des
systèmes.}

\begin{boxobjectifs}
\begin{itemize}
    \item Comprendre l'organisation d'un disque et d'une partition.
    \item Distinguer FAT, NTFS et EXT.
    \item Lire les métadonnées de fichiers (dates, propriétés).
    \item Repérer l'espace alloué et non alloué.
    \item Situer les artefacts à examiner.
\end{itemize}
\end{boxobjectifs}

\bigskip

% ============================================================================
\section*{Question directrice : où, dans l'image du disque, chercher la trace de la copie du fichier clients ?}
% ============================================================================

Le disque de M.~NDONGO garde, peut-être, la trace de la copie du fichier
clients sur la clé USB. Ce chapitre apprend à lire l'organisation du disque et
de ses systèmes de fichiers pour y trouver ces traces.

\medskip
\textbf{Plan du chapitre.}
\begin{enumerate}
    \item Disque et partitions (section~\ref{sec:c10-partitions}).
    \item FAT, NTFS, EXT (section~\ref{sec:c10-fs}).
    \item Les métadonnées de fichiers (section~\ref{sec:c10-meta}).
    \item Alloué et non alloué (section~\ref{sec:c10-alloue}).
\end{enumerate}

\separateur

% ============================================================================
\section{Disque et Partitions}
\label{sec:c10-partitions}
% ============================================================================

Un disque se divise en \termecle{partitions}, chacune formatée avec un
\termecle{système de fichiers} qui organise les données. La table des
partitions décrit cette organisation~; une partition cachée ou supprimée est,
en soi, un indice.

% ============================================================================
\section{FAT, NTFS, EXT}
\label{sec:c10-fs}
% ============================================================================

\begin{tableaucours}{L{2.4cm}L{4.4cm}L{5.4cm}}{%
    \tableauHeader{Système} &
    \tableauHeader{Contexte} &
    \tableauHeader{Intérêt forensique}
}
FAT / exFAT & Clés USB, cartes & Simple~; dates de création, modification, accès \\
NTFS & Windows & Riche~: journal, métadonnées, flux, horodatages multiples \\
EXT (2/3/4) & Linux & Inodes, horodatages, journal \\
\end{tableaucours}
\vspace{-0.8em}
\begin{center}{\small\itshape Tableau~10.1. Trois familles de systèmes de fichiers}\end{center}

% ============================================================================
\section{Les Métadonnées de Fichiers}
\label{sec:c10-meta}
% ============================================================================

Chaque fichier porte des \termecle{métadonnées}~: dates de création, de
modification, d'accès, de changement~; nom~; emplacement~; propriétaire. Ces
horodatages, croisés, aident à dater une action (copie, ouverture). On les
interprète avec prudence~: certaines opérations les modifient, d'autres non.

\begin{boxCRO}
\textbf{Lire sans sur-interpréter.} Un horodatage de modification prouve
qu'un fichier a changé, non \emph{qui} l'a changé ni \emph{pourquoi}~:
\cro{0}{+}{0} pour le fait, mais l'interprétation reste prudente. On
distingue le tracé (la date) du déduit (l'intention).
\end{boxCRO}

% ============================================================================
\section{Alloué et Non Alloué}
\label{sec:c10-alloue}
% ============================================================================

\begin{boxscenario}{Opération TSANAGA. La trace de la copie}
\textbf{Constatation.} Sur l'image NTFS du disque, le journal du système de
fichiers et les raccourcis récents montrent qu'un fichier
\texttt{clients.xlsx} a été ouvert puis copié vers un volume amovible le
3~janvier.

\medskip
\textbf{Analyse} (confiance élevée pour la trace, moyenne pour l'auteur). La
copie vers un amovible est établie par les artefacts NTFS. L'identité de la
personne devra être recoupée (session ouverte, autres indices). On date le
fait, on n'affirme pas encore l'auteur.
\end{boxscenario}

\begin{boxmetier}{Analyste forensique}
\textbf{Connaître le système de fichiers.} L'analyste adapte sa lecture au
système~: NTFS offre des horodatages multiples et un journal précieux~; FAT
en offre moins. Savoir ce que chaque système conserve évite les conclusions
erronées.
\end{boxmetier}

\separateur

\begin{boxresume}
\textbf{Ce qu'il faut retenir de ce chapitre~:}
\begin{itemize}
    \item \textbf{Un disque} se divise en partitions et systèmes de fichiers.
    \item \textbf{FAT, NTFS, EXT} conservent des traces de richesse inégale.
    \item \textbf{Les métadonnées} (horodatages) aident à dater une action.
    \item \textbf{On lit avec prudence}~: la date n'est pas l'auteur.
    \item \textbf{L'espace non alloué} garde ce qu'on croyait effacé.
\end{itemize}
\end{boxresume}

\begin{boxinfo}
\textbf{Points de contrôle de l'élève.} Je sais~: décrire disque et
partitions~; comparer FAT, NTFS, EXT~; lire des métadonnées~; distinguer
alloué et non alloué~; dater sans sur-interpréter.
\end{boxinfo}

\bigskip

\begin{boxexo}[title={\ding{45}\quad Exercice~10.1. Partition \niveaubase}]
Qu'est-ce qu'une partition~? Un système de fichiers~?
\end{boxexo}

\begin{boxexo}[title={\ding{45}\quad Exercice~10.2. NTFS \niveaubase}]
Pourquoi NTFS est-il plus riche que FAT pour l'enquête~?
\end{boxexo}

\begin{boxexo}[title={\ding{45}\quad Exercice~10.3. Métadonnées \niveaubase}]
Citez trois horodatages portés par un fichier.
\end{boxexo}

\begin{boxexo}[title={\ding{45}\quad Exercice~10.4. La date et l'auteur \niveaumoyen}]
Pourquoi un horodatage de modification ne désigne-t-il pas l'auteur~?
\end{boxexo}

\begin{boxexo}[title={\ding{45}\quad Exercice~10.5. TSANAGA \niveaumoyen}]
Quels artefacts NTFS montrent la copie de \texttt{clients.xlsx}~?
\end{boxexo}

\begin{boxexo}[title={\ding{45}\quad Exercice~10.6. Partition cachée \niveaumoyen}]
Pourquoi une partition cachée est-elle un indice~?
\end{boxexo}

\begin{boxexo}[title={\ding{45}\quad Exercice~10.7. CRO \niveaumoyen}]
Expliquez pourquoi on lit les horodatages avec prudence.
\end{boxexo}

\begin{boxexo}[title={\ding{45}\quad Exercice~10.8. EXT \niveaumoyen}]
Qu'examine-t-on dans un système de fichiers EXT~?
\end{boxexo}

\begin{boxexo}[title={\ding{45}\quad Exercice~10.9. Plan d'analyse \niveauexpert}]
Rédigez le plan d'analyse disque de l'image NTFS de TSANAGA.
\end{boxexo}

\begin{boxexo}[title={\ding{45}\quad Exercice~10.10. Constat \niveauexpert}]
Rédigez le constat de la copie de \texttt{clients.xlsx} vers l'amovible.
\end{boxexo}

\bigskip

\begin{boxinfo}
\textbf{Pour aller plus loin.}
\begin{itemize}
    \item ISO/IEC 27042~\cite{iso27042}~; NIST SP~800-86~\cite{nist80086}.
\end{itemize}
\end{boxinfo}

\begin{boxinfo}
\textbf{Transition.} Le disque garde aussi ce qu'on a voulu effacer~: le
Chapitre~\ref{chap:artefacts} traite des fichiers supprimés et des artefacts.
\end{boxinfo}
}{}
        \IfFileExists{partie4/P04_C11_supprimes_artefacts.tex}{% ============================================================================
% Fichier  : partie4/P04_C11_supprimes_artefacts.tex : FORN (LPI342)
% Sources : ISO/IEC 27042 ; NIST SP 800-86. Opération TSANAGA.
% ============================================================================

\chapter{Fichiers Supprimés, Fragments et Artefacts}
\label{chap:artefacts}
\chaptermark{Supprimés et artefacts}

\epigraph{%
    « Supprimer, c'est souvent seulement cacher. »%
}{Principe du cours}

\niveauChapitre{Intermédiaire à expert}{Chapitre~\ref{chap:disque}.}

\begin{boxobjectifs}
\begin{itemize}
    \item Comprendre pourquoi un fichier supprimé subsiste.
    \item Récupérer un fichier par ses fragments.
    \item Identifier un type de fichier par sa signature.
    \item Lire les artefacts d'usage (récents, raccourcis, corbeille).
    \item Qualifier un artefact récupéré.
\end{itemize}
\end{boxobjectifs}

\bigskip

% ============================================================================
\section*{Question directrice : le fichier clients, supprimé par M. NDONGO, peut-il encore être retrouvé ?}
% ============================================================================

M.~NDONGO a supprimé le fichier et vidé la corbeille. Pourtant, les données
peuvent subsister. Ce chapitre explique pourquoi, et comment les récupérer et
les qualifier.

\medskip
\textbf{Plan du chapitre.}
\begin{enumerate}
    \item Pourquoi un fichier subsiste (section~\ref{sec:c11-subsiste}).
    \item Récupérer par fragments (section~\ref{sec:c11-carving}).
    \item Les signatures de fichiers (section~\ref{sec:c11-signatures}).
    \item Les artefacts d'usage (section~\ref{sec:c11-usage}).
\end{enumerate}

\separateur

% ============================================================================
\section{Pourquoi un Fichier Subsiste}
\label{sec:c11-subsiste}
% ============================================================================

Supprimer un fichier ne l'efface pas~: le système marque son espace comme
réutilisable et retire son entrée de l'index, mais les données restent tant
qu'elles ne sont pas réécrites. D'où la possibilité de les
\termecle{récupérer}, en partie ou en totalité, jusqu'à leur écrasement.

% ============================================================================
\section{Récupérer par Fragments}
\label{sec:c11-carving}
% ============================================================================

\begin{boxdef}
\termecle{Récupération par signatures} (\textit{file carving})~: technique qui
reconstitue des fichiers à partir de leurs motifs de début et de fin dans
l'espace non alloué, sans s'appuyer sur l'index du système de fichiers. Elle
retrouve des fichiers dont l'entrée a disparu.
\end{boxdef}

% ============================================================================
\section{Les Signatures de Fichiers}
\label{sec:c11-signatures}
% ============================================================================

Un fichier commence souvent par une \termecle{signature} (quelques octets)
qui trahit son type, indépendamment de son extension. Un fichier renommé en
\texttt{.txt} mais commençant par la signature d'une archive reste une
archive~: la signature prime sur l'extension.

\begin{boxalert}
\textbf{Violation typique.} Se fier à l'extension d'un fichier.

\smallskip
\textbf{Règle~:} l'extension se change en un clic. On identifie le type par la
signature, pas par le nom~: c'est une parade classique de dissimulation.
\end{boxalert}

% ============================================================================
\section{Les Artefacts d'Usage}
\label{sec:c11-usage}
% ============================================================================

\begin{tableaucours}{L{3.6cm}L{8.6cm}}{%
    \tableauHeader{Artefact} &
    \tableauHeader{Ce qu'il révèle}
}
Fichiers récents, raccourcis & Documents ouverts, chemins, volumes amovibles \\
Corbeille & Fichiers supprimés récupérables, date \\
Journaux d'événements & Branchement de périphériques, sessions \\
Fichiers temporaires, caches & Copies de travail, miniatures \\
\end{tableaucours}
\vspace{-0.8em}
\begin{center}{\small\itshape Tableau~11.1. Artefacts d'usage utiles}\end{center}

\begin{boxscenario}{Opération TSANAGA. Le fichier retrouvé}
\textbf{Constatation.} Dans l'espace non alloué, la récupération par
signatures reconstitue un tableur dont la signature est celle d'un classeur,
bien que l'entrée ait disparu. Son contenu correspond au fichier clients. Un
raccourci récent pointe vers un volume amovible, le 3~janvier.

\medskip
\textbf{Analyse} (confiance élevée pour l'existence, moyenne pour l'auteur).
Le fichier supprimé est récupéré et son contenu établi~; la copie vers
l'amovible est corroborée par l'artefact de raccourci. Le faisceau se
renforce.
\end{boxscenario}

\begin{boxobjet}{Fichier récupéré : clients (classeur)}
\textbf{Nature~:} classeur récupéré par signatures dans l'espace non alloué.
\textbf{Source~:} image disque TSA-2026-01 (condensat au dossier).
\textbf{Pertinence~:} contenu correspondant au fichier clients.
\textbf{Fiabilité~:} intègre (travail sur l'image). \textbf{Opposabilité~:}
chaîne de possession continue.
\end{boxobjet}

\begin{boxmetier}{Analyste forensique}
\textbf{Ne pas se fier au nom.} L'analyste identifie les fichiers par leur
signature, fouille l'espace non alloué, et recoupe les artefacts d'usage. Ce
que l'utilisateur croyait caché resurgit souvent.
\end{boxmetier}

\separateur

\begin{boxresume}
\textbf{Ce qu'il faut retenir de ce chapitre~:}
\begin{itemize}
    \item \textbf{Supprimer n'est pas effacer}~: récupération possible.
    \item \textbf{La récupération par signatures} reconstitue sans l'index.
    \item \textbf{La signature} prime sur l'extension.
    \item \textbf{Les artefacts d'usage} révèlent documents et périphériques.
    \item \textbf{On qualifie} un fichier récupéré.
\end{itemize}
\end{boxresume}

\begin{boxinfo}
\textbf{Points de contrôle de l'élève.} Je sais~: expliquer la subsistance
d'un fichier supprimé~; décrire la récupération par signatures~; identifier un
type par signature~; lire les artefacts d'usage~; qualifier un fichier
récupéré.
\end{boxinfo}

\bigskip

\begin{boxexo}[title={\ding{45}\quad Exercice~11.1. Subsister \niveaubase}]
Pourquoi un fichier supprimé peut-il être récupéré~?
\end{boxexo}

\begin{boxexo}[title={\ding{45}\quad Exercice~11.2. Carving \niveaubase}]
Qu'est-ce que la récupération par signatures~?
\end{boxexo}

\begin{boxexo}[title={\ding{45}\quad Exercice~11.3. Extension \niveaubase}]
Pourquoi ne pas se fier à l'extension d'un fichier~?
\end{boxexo}

\begin{boxexo}[title={\ding{45}\quad Exercice~11.4. Signature \niveaumoyen}]
Comment identifier un fichier renommé pour dissimuler son type~?
\end{boxexo}

\begin{boxexo}[title={\ding{45}\quad Exercice~11.5. TSANAGA \niveaumoyen}]
Comment le fichier clients supprimé a-t-il été retrouvé et corroboré~?
\end{boxexo}

\begin{boxexo}[title={\ding{45}\quad Exercice~11.6. Raccourci \niveaumoyen}]
Que révèle un raccourci récent pointant vers un volume amovible~?
\end{boxexo}

\begin{boxexo}[title={\ding{45}\quad Exercice~11.7. Écrasement \niveaumoyen}]
Qu'est-ce qui limite la récupération d'un fichier supprimé~?
\end{boxexo}

\begin{boxexo}[title={\ding{45}\quad Exercice~11.8. Corbeille \niveaumoyen}]
Que peut révéler la corbeille, même vidée~?
\end{boxexo}

\begin{boxexo}[title={\ding{45}\quad Exercice~11.9. Faisceau \niveauexpert}]
Construisez le faisceau d'indices de la copie du fichier clients.
\end{boxexo}

\begin{boxexo}[title={\ding{45}\quad Exercice~11.10. Qualification \niveauexpert}]
Qualifiez complètement le fichier récupéré (cinq rubriques).
\end{boxexo}

\bigskip

\begin{boxinfo}
\textbf{Pour aller plus loin.}
\begin{itemize}
    \item ISO/IEC 27042~\cite{iso27042}~; NIST SP~800-86~\cite{nist80086}.
\end{itemize}
\end{boxinfo}

\begin{boxinfo}
\textbf{Transition.} Le disque ne dit pas tout~: la mémoire garde l'état
vivant. Le Chapitre~\ref{chap:memoire} traite du forensique mémoire.
\end{boxinfo}
}{}
        \IfFileExists{partie4/P04_C12_forensique_memoire.tex}{% ============================================================================
% Fichier  : partie4/P04_C12_forensique_memoire.tex : FORN (LPI342)
% Sources : NIST SP 800-86 ; RFC 3227. Opération TSANAGA.
% ============================================================================

\chapter{Forensique Mémoire}
\label{chap:memoire}
\chaptermark{Forensique mémoire}

\epigraph{%
    « La mémoire vive avoue ce que le disque dissimule. »%
}{Principe du cours}

\niveauChapitre{Expert}{Chapitres~\ref{chap:acquisition} et~\ref{chap:disque}.}

\begin{boxobjectifs}
\begin{itemize}
    \item Comprendre ce que révèle la mémoire vive.
    \item Analyser une capture mémoire.
    \item Repérer processus, connexions, secrets.
    \item Relier la mémoire au disque et au réseau.
    \item Qualifier un élément de mémoire.
\end{itemize}
\end{boxobjectifs}

\bigskip

% ============================================================================
\section*{Question directrice : que contenait la mémoire de l'ordinateur de M. NDONGO au moment de la saisie ?}
% ============================================================================

La capture mémoire réalisée à chaud (Chapitre~\ref{chap:acquisition}) contient
l'état vivant du système~: c'est là que se trouvait la clé du conteneur
chiffré. Ce chapitre apprend à l'analyser.

\medskip
\textbf{Plan du chapitre.}
\begin{enumerate}
    \item Ce que révèle la mémoire (section~\ref{sec:c12-revele}).
    \item Analyser une capture (section~\ref{sec:c12-analyser}).
    \item Secrets et chiffrement (section~\ref{sec:c12-secrets}).
    \item Recouper avec le disque (section~\ref{sec:c12-recouper}).
\end{enumerate}

\separateur

% ============================================================================
\section{Ce que Révèle la Mémoire}
\label{sec:c12-revele}
% ============================================================================

\begin{tableaucours}{L{4.0cm}L{8.2cm}}{%
    \tableauHeader{Élément} &
    \tableauHeader{Intérêt}
}
Processus en cours & Programmes actifs, parent suspect \\
Connexions réseau & Destinations contactées au moment de la saisie \\
Secrets déchiffrés & Clés, mots de passe, jetons en clair \\
Fichiers ouverts & Documents en cours d'usage \\
Fragments & Textes, adresses, en mémoire mais non sur disque \\
\end{tableaucours}
\vspace{-0.8em}
\begin{center}{\small\itshape Tableau~12.1. Ce que révèle la mémoire vive}\end{center}

% ============================================================================
\section{Analyser une Capture}
\label{sec:c12-analyser}
% ============================================================================

\begin{boxprocedure}
\textbf{Analyser une capture mémoire}
\begin{enumerate}
    \item Vérifier l'intégrité de la capture (condensat).
    \item Lister les processus et leurs liens de parenté.
    \item Lister les connexions réseau ouvertes.
    \item Rechercher secrets, clés, fragments pertinents.
    \item Recouper avec le disque et le réseau.
\end{enumerate}
\end{boxprocedure}

% ============================================================================
\section{Secrets et Chiffrement}
\label{sec:c12-secrets}
% ============================================================================

\begin{boxscenario}{Opération TSANAGA. La clé en mémoire}
\textbf{Constatation.} La capture mémoire contient la clé du conteneur
chiffré, en clair (le conteneur était monté à la saisie). Elle permet d'ouvrir
le conteneur et d'en analyser le contenu.

\medskip
\textbf{Analyse} (confiance élevée). Sans la capture à chaud, le conteneur
serait resté inaccessible. La clé récupérée a permis d'établir que le
conteneur contenait une copie du fichier clients. L'arbitrage de l'acquisition
à chaud trouve ici sa justification.
\end{boxscenario}

\begin{boxCRO}
\textbf{Une clé en mémoire, un pouvoir à encadrer.} Déchiffrer un conteneur
grâce à une clé trouvée en mémoire reste dans le périmètre de la réquisition~:
\cro{0}{+}{+}. En sortir (déchiffrer d'autres supports non visés) violerait la
confidentialité et le mandat.
\end{boxCRO}

% ============================================================================
\section{Recouper avec le Disque}
\label{sec:c12-recouper}
% ============================================================================

La mémoire, seule, est un instantané~; recoupée avec le disque (le fichier
récupéré) et le réseau (les destinations), elle devient un faisceau. La force
de l'enquête tient à la convergence des domaines sur une même histoire.

\begin{boxmetier}{Analyste forensique}
\textbf{La mémoire, fenêtre sur l'instant.} L'analyste traite la capture
mémoire comme une photographie de l'instant de la saisie~: précieuse, unique,
à recouper. Il en extrait ce qui sert l'affaire, dans le périmètre.
\end{boxmetier}

\separateur

\begin{boxresume}
\textbf{Ce qu'il faut retenir de ce chapitre~:}
\begin{itemize}
    \item \textbf{La mémoire} révèle processus, connexions, secrets,
          fragments.
    \item \textbf{On vérifie} l'intégrité de la capture avant de l'analyser.
    \item \textbf{Les secrets en clair} (clés) ouvrent les conteneurs
          chiffrés.
    \item \textbf{On reste dans le périmètre} de la réquisition.
    \item \textbf{On recoupe} mémoire, disque, réseau.
\end{itemize}
\end{boxresume}

\begin{boxinfo}
\textbf{Points de contrôle de l'élève.} Je sais~: citer ce que révèle la
mémoire~; analyser une capture~; exploiter un secret dans le périmètre~;
recouper les domaines~; qualifier un élément de mémoire.
\end{boxinfo}

\bigskip

\begin{boxexo}[title={\ding{45}\quad Exercice~12.1. Révélations \niveaubase}]
Citez trois éléments que révèle la mémoire vive.
\end{boxexo}

\begin{boxexo}[title={\ding{45}\quad Exercice~12.2. Intégrité \niveaubase}]
Que vérifie-t-on avant d'analyser une capture mémoire~?
\end{boxexo}

\begin{boxexo}[title={\ding{45}\quad Exercice~12.3. La clé \niveaubase}]
Qu'a permis la clé trouvée en mémoire dans TSANAGA~?
\end{boxexo}

\begin{boxexo}[title={\ding{45}\quad Exercice~12.4. Périmètre \niveaumoyen}]
Peut-on utiliser une clé trouvée en mémoire pour déchiffrer un support non
visé~? Pourquoi~?
\end{boxexo}

\begin{boxexo}[title={\ding{45}\quad Exercice~12.5. Instantané \niveaumoyen}]
Pourquoi la mémoire est-elle un instantané à recouper~?
\end{boxexo}

\begin{boxexo}[title={\ding{45}\quad Exercice~12.6. CRO \niveaumoyen}]
Expliquez l'analyse CRO de l'usage d'une clé trouvée en mémoire.
\end{boxexo}

\begin{boxexo}[title={\ding{45}\quad Exercice~12.7. Parenté \niveaumoyen}]
Pourquoi examiner les liens de parenté entre processus~?
\end{boxexo}

\begin{boxexo}[title={\ding{45}\quad Exercice~12.8. Convergence \niveaumoyen}]
Comment la mémoire se recoupe-t-elle avec le disque dans TSANAGA~?
\end{boxexo}

\begin{boxexo}[title={\ding{45}\quad Exercice~12.9. Plan d'analyse \niveauexpert}]
Rédigez le plan d'analyse de la capture mémoire de TSANAGA.
\end{boxexo}

\begin{boxexo}[title={\ding{45}\quad Exercice~12.10. Justification \niveauexpert}]
Rédigez la note qui justifie l'exploitation de la clé trouvée en mémoire.
\end{boxexo}

\bigskip

\begin{boxinfo}
\textbf{Pour aller plus loin.}
\begin{itemize}
    \item NIST SP~800-86~\cite{nist80086}~; RFC~3227~\cite{rfc3227}.
\end{itemize}
\end{boxinfo}

\begin{boxinfo}
\textbf{Transition.} Reste la trace des échanges~: le
Chapitre~\ref{chap:reseau} traite du forensique réseau.
\end{boxinfo}
}{}
        \IfFileExists{partie4/P04_C13_forensique_reseau.tex}{% ============================================================================
% Fichier  : partie4/P04_C13_forensique_reseau.tex : FORN (LPI342)
% Sources : RFC 3227 ; NIST SP 800-86 ; MITRE ATT&CK. Opération TSANAGA.
% Le bureau d'analyse travaille sur des journaux de flux exportés (pas de
% capture graphique).
% ============================================================================

\chapter{Forensique Réseau}
\label{chap:reseau}
\chaptermark{Forensique réseau}

\epigraph{%
    « Le disque dit ce qu'on a gardé~; le réseau dit ce qu'on a envoyé. »%
}{Principe du cours}

\niveauChapitre{Intermédiaire à expert}{Chapitre~\ref{chap:memoire}~; notions
de réseaux.}

\begin{boxobjectifs}
\begin{itemize}
    \item Identifier les sources de traces réseau.
    \item Repérer une exfiltration et un canal de pilotage.
    \item Lire des journaux de flux exportés.
    \item Analyser les échanges suspects (pilotage, dissimulation par DNS).
    \item Reconstituer une chronologie réseau.
\end{itemize}
\end{boxobjectifs}

\bigskip

% ============================================================================
\section*{Question directrice : le fichier clients a-t-il quitté le réseau de LiZbiZ, et vers où ?}
% ============================================================================

Au-delà de la copie sur clé, M.~NDONGO a-t-il aussi envoyé le fichier par le
réseau~? Les journaux de flux et de mandataire gardent la trace des envois.
Ce chapitre apprend à les lire.

\medskip
\textbf{Plan du chapitre.}
\begin{enumerate}
    \item Les sources réseau (section~\ref{sec:c13-sources}).
    \item Exfiltration et pilotage (section~\ref{sec:c13-exfil}).
    \item Les échanges suspects (section~\ref{sec:c13-suspects}).
    \item La chronologie réseau (section~\ref{sec:c13-chrono}).
\end{enumerate}

\separateur

% ============================================================================
\section{Les Sources Réseau}
\label{sec:c13-sources}
% ============================================================================

\begin{tableaucours}{L{3.6cm}L{8.6cm}}{%
    \tableauHeader{Source} &
    \tableauHeader{Ce qu'elle révèle}
}
Journaux de flux (pare-feu) & Qui parle à qui, quand, quel volume \\
Journaux du mandataire web & Destinations web, volumes, pièces envoyées \\
Journaux DNS & Noms résolus, dissimulation par DNS \\
Journaux de messagerie & Courriels, pièces jointes, destinataires \\
\end{tableaucours}
\vspace{-0.8em}
\begin{center}{\small\itshape Tableau~13.1. Sources de traces réseau}\end{center}

\begin{boxalert}
\textbf{Contrainte du cours.} Le bureau d'analyse de NDAP SUKULU n'a pas
d'outil de capture graphique.

\smallskip
\textbf{Règle~:} on travaille sur des \emph{journaux de flux exportés} (au
format tableur), suffisants pour établir les faits et parfaitement opposables.
La méthode prime sur l'outil.
\end{boxalert}

% ============================================================================
\section{Exfiltration et Pilotage}
\label{sec:c13-exfil}
% ============================================================================

Deux signatures réseau comptent~: l'\termecle{exfiltration} (un envoi sortant
anormal vers une destination externe) et, dans les affaires d'intrusion, le
\termecle{canal de pilotage} (petites connexions régulières commandant un code
malveillant). Dans TSANAGA, on cherche surtout l'envoi du fichier clients.

% ============================================================================
\section{Les Échanges Suspects}
\label{sec:c13-suspects}
% ============================================================================

Certains envois se dissimulent~: données glissées dans des requêtes DNS,
transferts vers un service d'hébergement personnel, courriels à une adresse
externe. L'analyste repère les destinations inhabituelles et les volumes qui
ne correspondent pas à l'usage normal.

% ============================================================================
\section{La Chronologie Réseau}
\label{sec:c13-chrono}
% ============================================================================

\begin{boxscenario}{Opération TSANAGA. L'envoi du 5 janvier}
\textbf{Constatation.} Les journaux du mandataire montrent, le 5~janvier, un
envoi de 18~Mo depuis le poste de M.~NDONGO vers un service d'hébergement
personnel, deux jours après la copie sur clé.

\medskip
\textbf{Analyse} (confiance élevée pour l'envoi, moyenne pour le contenu).
L'envoi est établi~; que son contenu soit le fichier clients est fortement
probable (taille, date), mais se recoupe avec le disque. On date le fait, on
qualifie la probabilité.
\end{boxscenario}

\begin{boxmetier}{Analyste forensique}
\textbf{Suivre la donnée.} L'analyste trace le chemin de la donnée~: copiée
sur une clé, envoyée par le réseau. Chaque canal laisse une trace~; les
recouper reconstitue l'histoire complète de l'exfiltration.
\end{boxmetier}

\separateur

\begin{boxresume}
\textbf{Ce qu'il faut retenir de ce chapitre~:}
\begin{itemize}
    \item \textbf{Flux, mandataire, DNS, messagerie}~: les sources réseau.
    \item \textbf{Exfiltration} (envoi sortant) et \textbf{pilotage} (petites
          connexions régulières).
    \item \textbf{Des journaux de flux exportés} suffisent.
    \item \textbf{Les échanges se dissimulent}~: destinations inhabituelles.
    \item \textbf{On reconstitue} la chronologie et on recoupe.
\end{itemize}
\end{boxresume}

\begin{boxinfo}
\textbf{Points de contrôle de l'élève.} Je sais~: citer les sources réseau~;
reconnaître exfiltration et pilotage~; lire des flux exportés~; repérer un
échange suspect~; dater un envoi.
\end{boxinfo}

\bigskip

\begin{boxexo}[title={\ding{45}\quad Exercice~13.1. Sources \niveaubase}]
Quelle source révèle les pièces jointes envoyées par courriel~?
\end{boxexo}

\begin{boxexo}[title={\ding{45}\quad Exercice~13.2. Exfiltration \niveaubase}]
Quelle signature réseau trahit un envoi sortant anormal~?
\end{boxexo}

\begin{boxexo}[title={\ding{45}\quad Exercice~13.3. Sans capture \niveaubase}]
Pourquoi des journaux de flux exportés suffisent-ils~?
\end{boxexo}

\begin{boxexo}[title={\ding{45}\quad Exercice~13.4. Dissimulation \niveaumoyen}]
Comment des données peuvent-elles se dissimuler dans des requêtes DNS~?
\end{boxexo}

\begin{boxexo}[title={\ding{45}\quad Exercice~13.5. TSANAGA \niveaumoyen}]
Qu'établissent les journaux du mandataire le 5~janvier~?
\end{boxexo}

\begin{boxexo}[title={\ding{45}\quad Exercice~13.6. Contenu \niveaumoyen}]
Pourquoi l'envoi est-il « établi » mais son contenu « probable »~?
\end{boxexo}

\begin{boxexo}[title={\ding{45}\quad Exercice~13.7. Recouper \niveaumoyen}]
Comment recouper l'envoi réseau avec le disque~?
\end{boxexo}

\begin{boxexo}[title={\ding{45}\quad Exercice~13.8. Destination \niveaumoyen}]
Comment repère-t-on une destination inhabituelle dans les flux~?
\end{boxexo}

\begin{boxexo}[title={\ding{45}\quad Exercice~13.9. Chronologie \niveauexpert}]
Dressez la chronologie réseau de l'exfiltration dans TSANAGA.
\end{boxexo}

\begin{boxexo}[title={\ding{45}\quad Exercice~13.10. Chemin de la donnée \niveauexpert}]
Reconstituez le chemin complet du fichier clients (clé et réseau).
\end{boxexo}

\bigskip

\begin{boxinfo}
\textbf{Pour aller plus loin.}
\begin{itemize}
    \item RFC~3227~\cite{rfc3227}~; NIST SP~800-86~\cite{nist80086}~;
          MITRE ATT\&CK~\cite{mitreattack}.
\end{itemize}
\end{boxinfo}

\begin{boxinfo}
\textbf{Transition.} L'enquête se déplace vers la poche du suspect~: le
Chapitre~\ref{chap:mobile} traite du forensique mobile.
\end{boxinfo}
}{}
        \IfFileExists{partie4/P04_C14_forensique_mobile.tex}{% ============================================================================
% Fichier  : partie4/P04_C14_forensique_mobile.tex : FORN (LPI342)
% Sources : NIST SP 800-101r1 ; ISO/IEC 27037. Opération TSANAGA.
% ============================================================================

\chapter{Forensique Mobile}
\label{chap:mobile}
\chaptermark{Forensique mobile}

\epigraph{%
    « Le téléphone est le témoin qui ne quitte jamais son propriétaire. »%
}{Principe du cours}

\niveauChapitre{Intermédiaire à expert}{Chapitre~\ref{chap:disque}~; notions
de systèmes mobiles.}

\begin{boxobjectifs}
\begin{itemize}
    \item Comprendre les spécificités de l'enquête sur smartphone.
    \item Identifier les traces (messages, applications, géolocalisation).
    \item Connaître les niveaux d'extraction.
    \item Respecter le cadre juridique du mobile.
    \item Qualifier un artefact mobile.
\end{itemize}
\end{boxobjectifs}

\bigskip

% ============================================================================
\section*{Question directrice : le téléphone de M. NDONGO confirme-t-il la revente des données ?}
% ============================================================================

Le téléphone saisi peut contenir messages, contacts, et traces de
négociation. L'enquête mobile obéit à des règles propres. Ce chapitre les
présente.

\medskip
\textbf{Plan du chapitre.}
\begin{enumerate}
    \item Les spécificités du mobile (section~\ref{sec:c14-spec}).
    \item Les traces mobiles (section~\ref{sec:c14-traces}).
    \item Les niveaux d'extraction (section~\ref{sec:c14-extraction}).
    \item Le cadre juridique (section~\ref{sec:c14-droit}).
\end{enumerate}

\separateur

% ============================================================================
\section{Les Spécificités du Mobile}
\label{sec:c14-spec}
% ============================================================================

Un smartphone diffère d'un ordinateur~: stockage chiffré par défaut,
verrouillage, connexion permanente (risque d'effacement à distance),
diversité des systèmes (Android, iOS). On l'isole des réseaux dès la saisie
(mode avion, sachet de blindage) pour éviter l'effacement.

% ============================================================================
\section{Les Traces Mobiles}
\label{sec:c14-traces}
% ============================================================================

\begin{tableaucours}{L{3.6cm}L{8.6cm}}{%
    \tableauHeader{Trace} &
    \tableauHeader{Ce qu'elle révèle}
}
Messages & SMS, messageries, conversations \\
Journaux d'appels & Qui, quand, combien de temps \\
Applications & Données locales, historiques, pièces jointes \\
Géolocalisation & Lieux fréquentés, déplacements \\
Métadonnées de photos & Date, lieu de prise de vue \\
\end{tableaucours}
\vspace{-0.8em}
\begin{center}{\small\itshape Tableau~14.1. Traces d'un smartphone~\cite{nist800101}}\end{center}

% ============================================================================
\section{Les Niveaux d'Extraction}
\label{sec:c14-extraction}
% ============================================================================

\begin{tableaucours}{L{3.6cm}L{8.6cm}}{%
    \tableauHeader{Niveau} &
    \tableauHeader{Portée}
}
Manuel & Lecture à l'écran, photographiée~; limité, mais simple \\
Logique & Données accessibles par l'interface (messages, contacts) \\
Physique & Copie bas niveau~; plus complet, plus complexe \\
\end{tableaucours}
\vspace{-0.8em}
\begin{center}{\small\itshape Tableau~14.2. Niveaux d'extraction mobile}\end{center}

% ============================================================================
\section{Le Cadre Juridique}
\label{sec:c14-droit}
% ============================================================================

\begin{boxCRO}
\textbf{Le mobile, intime par nature.} Un téléphone contient la vie privée
entière de son propriétaire~: messages intimes, santé, famille. L'examiner
sans limite dégrade la \cro{-}{0}{0}. On s'en tient strictement au périmètre
de la réquisition et l'on protège les tiers (correspondants).
\end{boxCRO}

\begin{boxscenario}{Opération TSANAGA. Les messages de négociation}
\textbf{Constatation.} L'extraction logique du téléphone, dans le périmètre,
révèle des messages datés de janvier évoquant la remise d'un « fichier » contre
paiement, à un contact extérieur.

\medskip
\textbf{Analyse} (confiance élevée pour les messages, moyenne pour leur
portée exacte). Les messages corroborent fortement la revente~; leur
interprétation précise (quel fichier, quel montant) se recoupe avec le disque
et le réseau. On rapporte les messages, à charge et à décharge.
\end{boxscenario}

\begin{boxmetier}{Enquêteur en cybercriminalité et expert judiciaire}
\textbf{Isoler, puis extraire dans le périmètre.} L'enquêteur isole le
téléphone des réseaux dès la saisie, choisit le niveau d'extraction adapté, et
limite l'examen au périmètre. Le respect de l'intimité conditionne la
recevabilité.
\end{boxmetier}

\separateur

\begin{boxresume}
\textbf{Ce qu'il faut retenir de ce chapitre~:}
\begin{itemize}
    \item \textbf{Le mobile} est chiffré, verrouillé, connecté~: on l'isole.
    \item \textbf{Les traces}~: messages, appels, applications,
          géolocalisation.
    \item \textbf{Trois niveaux d'extraction}~: manuel, logique, physique.
    \item \textbf{Le mobile est intime}~: on s'en tient au périmètre.
    \item \textbf{On recoupe} avec le disque et le réseau.
\end{itemize}
\end{boxresume}

\begin{boxinfo}
\textbf{Points de contrôle de l'élève.} Je sais~: expliquer les spécificités
du mobile~; citer les traces~; distinguer les niveaux d'extraction~; respecter
le périmètre~; qualifier un artefact mobile.
\end{boxinfo}

\bigskip

\begin{boxexo}[title={\ding{45}\quad Exercice~14.1. Spécificités \niveaubase}]
Pourquoi isole-t-on un téléphone des réseaux dès la saisie~?
\end{boxexo}

\begin{boxexo}[title={\ding{45}\quad Exercice~14.2. Traces \niveaubase}]
Citez quatre traces d'un smartphone.
\end{boxexo}

\begin{boxexo}[title={\ding{45}\quad Exercice~14.3. Niveaux \niveaubase}]
Distinguez extraction logique et physique.
\end{boxexo}

\begin{boxexo}[title={\ding{45}\quad Exercice~14.4. Effacement \niveaumoyen}]
Quel risque court un téléphone resté connecté après la saisie~?
\end{boxexo}

\begin{boxexo}[title={\ding{45}\quad Exercice~14.5. Intimité \niveaumoyen}]
Pourquoi le mobile exige-t-il une attention particulière à l'intimité~?
\end{boxexo}

\begin{boxexo}[title={\ding{45}\quad Exercice~14.6. TSANAGA \niveaumoyen}]
Que révèle l'extraction logique du téléphone de M.~NDONGO~?
\end{boxexo}

\begin{boxexo}[title={\ding{45}\quad Exercice~14.7. Géolocalisation \niveaumoyen}]
Que peut apporter la géolocalisation à l'enquête~?
\end{boxexo}

\begin{boxexo}[title={\ding{45}\quad Exercice~14.8. Tiers \niveaumoyen}]
Comment protéger les correspondants dont les messages figurent sur le
téléphone~?
\end{boxexo}

\begin{boxexo}[title={\ding{45}\quad Exercice~14.9. Procédure \niveauexpert}]
Rédigez la procédure de saisie et d'extraction mobile pour TSANAGA.
\end{boxexo}

\begin{boxexo}[title={\ding{45}\quad Exercice~14.10. Faisceau \niveauexpert}]
Intégrez les messages mobiles au faisceau d'indices de l'affaire.
\end{boxexo}

\bigskip

\begin{boxinfo}
\textbf{Pour aller plus loin.}
\begin{itemize}
    \item NIST SP~800-101r1~\cite{nist800101}~; ISO/IEC 27037~\cite{iso27037}.
\end{itemize}
\end{boxinfo}

\begin{boxinfo}
\textbf{Transition.} Les données vivent aussi hors de l'appareil~: le
Chapitre~\ref{chap:nuage} traite du forensique du nuage et des applications.
\end{boxinfo}
}{}
        \IfFileExists{partie4/P04_C15_forensique_nuage.tex}{% ============================================================================
% Fichier  : partie4/P04_C15_forensique_nuage.tex : FORN (LPI342)
% Sources : ISO/IEC 27037, 27042 ; NIST (cloud). Opération TSANAGA.
% ============================================================================

\chapter{Forensique du Nuage et des Applications}
\label{chap:nuage}
\chaptermark{Forensique du nuage}

\epigraph{%
    « Dans le nuage, la preuve n'est pas chez vous~;\\
    il faut la demander, dans les formes. »%
}{Principe du cours}

\niveauChapitre{Expert}{Chapitre~\ref{chap:mobile}~; notions de services en
ligne.}

\begin{boxobjectifs}
\begin{itemize}
    \item Comprendre les défis du forensique dans le nuage.
    \item Identifier les traces côté client et côté service.
    \item Obtenir une preuve auprès d'un hébergeur, dans les formes.
    \item Analyser les artefacts d'applications (réseaux sociaux, messageries).
    \item Respecter le droit et la coopération internationale.
\end{itemize}
\end{boxobjectifs}

\bigskip

% ============================================================================
\section*{Question directrice : le fichier a été envoyé vers un service d'hébergement. Comment en obtenir la preuve ?}
% ============================================================================

L'envoi du 5~janvier (Chapitre~\ref{chap:reseau}) pointe vers un service
d'hébergement personnel, hébergé à l'étranger. La preuve n'est plus sur un
support saisi, mais chez un tiers. Ce chapitre traite de ce cas.

\medskip
\textbf{Plan du chapitre.}
\begin{enumerate}
    \item Les défis du nuage (section~\ref{sec:c15-defis}).
    \item Côté client, côté service (section~\ref{sec:c15-cotes}).
    \item Obtenir dans les formes (section~\ref{sec:c15-obtenir}).
    \item Les applications (section~\ref{sec:c15-apps}).
\end{enumerate}

\separateur

% ============================================================================
\section{Les Défis du Nuage}
\label{sec:c15-defis}
% ============================================================================

Dans le nuage, la donnée est répartie, chez un tiers, souvent à l'étranger.
L'enquêteur ne peut ni la saisir physiquement, ni l'imager comme un disque. Il
doit la \emph{demander} au service, dans les formes légales, et composer avec
la localisation des données et la coopération internationale.

% ============================================================================
\section{Côté Client, Côté Service}
\label{sec:c15-cotes}
% ============================================================================

\begin{tableaucours}{L{3.6cm}L{8.6cm}}{%
    \tableauHeader{Côté} &
    \tableauHeader{Traces accessibles}
}
Client (le poste) & Fichiers de synchronisation, caches, jetons, historiques \\
Service (l'hébergeur) & Comptes, journaux d'accès, fichiers, sauvegardes \\
\end{tableaucours}
\vspace{-0.8em}
\begin{center}{\small\itshape Tableau~15.1. Deux côtés de la preuve en nuage}\end{center}

Souvent, le poste saisi garde déjà des traces du service (dossier de
synchronisation, jeton d'accès), exploitables sans attendre le tiers.

% ============================================================================
\section{Obtenir dans les Formes}
\label{sec:c15-obtenir}
% ============================================================================

\begin{boxjuris}
Obtenir une preuve auprès d'un hébergeur étranger passe par les voies
légales~: réquisition, coopération judiciaire internationale (Convention de
Budapest), demande au point de contact du service. Accéder directement au
compte, même avec un jeton trouvé, hors de ce cadre, exposerait à
l'irrecevabilité et à l'illégalité.
\end{boxjuris}

\begin{boxalert}
\textbf{Violation typique.} Se connecter au service avec un jeton d'accès
trouvé sur le poste, pour « aller voir ».

\smallskip
\textbf{Règle~:} un jeton trouvé établit un lien~; il ne vaut pas autorisation
d'accéder au service. On documente le jeton comme trace locale, et l'on
obtient le contenu du service par la voie légale.
\end{boxalert}

% ============================================================================
\section{Les Applications}
\label{sec:c15-apps}
% ============================================================================

\begin{boxscenario}{Opération TSANAGA. La trace de synchronisation}
\textbf{Constatation.} Sur le poste saisi, le dossier de synchronisation du
service d'hébergement contient une copie locale du fichier clients et un
journal de téléversement daté du 5~janvier.

\medskip
\textbf{Analyse} (confiance élevée). La trace locale, recevable (poste saisi,
chaîne continue), établit le téléversement sans attendre l'hébergeur. La
demande au service est engagée en parallèle, dans les formes, pour corroborer.
\end{boxscenario}

\begin{boxmetier}{Analyste forensique et juriste TIC}
\textbf{Chercher d'abord côté client.} Avant d'engager une longue coopération
internationale, l'analyste exploite les traces du service déjà présentes sur
le poste saisi~: elles sont souvent suffisantes et immédiatement recevables.
\end{boxmetier}

\separateur

\begin{boxresume}
\textbf{Ce qu'il faut retenir de ce chapitre~:}
\begin{itemize}
    \item \textbf{Le nuage}~: donnée répartie, chez un tiers, souvent à
          l'étranger.
    \item \textbf{Côté client}, le poste garde souvent des traces du service.
    \item \textbf{Côté service}, on obtient dans les formes (réquisition,
          Budapest).
    \item \textbf{Un jeton trouvé} n'autorise pas l'accès au service.
    \item \textbf{On exploite d'abord les traces locales}, recevables.
\end{itemize}
\end{boxresume}

\begin{boxinfo}
\textbf{Points de contrôle de l'élève.} Je sais~: expliquer les défis du
nuage~; distinguer traces client et service~; obtenir une preuve dans les
formes~; exploiter une trace locale~; respecter le droit.
\end{boxinfo}

\bigskip

\begin{boxexo}[title={\ding{45}\quad Exercice~15.1. Défi \niveaubase}]
Pourquoi ne peut-on pas imager une donnée dans le nuage comme un disque~?
\end{boxexo}

\begin{boxexo}[title={\ding{45}\quad Exercice~15.2. Deux côtés \niveaubase}]
Quelles traces trouve-t-on côté client~? Côté service~?
\end{boxexo}

\begin{boxexo}[title={\ding{45}\quad Exercice~15.3. Jeton \niveaubase}]
Un jeton trouvé autorise-t-il à se connecter au service~?
\end{boxexo}

\begin{boxexo}[title={\ding{45}\quad Exercice~15.4. Dans les formes \niveaumoyen}]
Par quelles voies obtient-on une preuve auprès d'un hébergeur étranger~?
\end{boxexo}

\begin{boxexo}[title={\ding{45}\quad Exercice~15.5. TSANAGA \niveaumoyen}]
Quelle trace locale établit le téléversement, sans attendre l'hébergeur~?
\end{boxexo}

\begin{boxexo}[title={\ding{45}\quad Exercice~15.6. Budapest \niveaumoyen}]
Quel rôle joue la Convention de Budapest ici~?
\end{boxexo}

\begin{boxexo}[title={\ding{45}\quad Exercice~15.7. Synchronisation \niveaumoyen}]
Pourquoi un dossier de synchronisation est-il précieux~?
\end{boxexo}

\begin{boxexo}[title={\ding{45}\quad Exercice~15.8. Priorité \niveaumoyen}]
Pourquoi exploiter d'abord les traces locales~?
\end{boxexo}

\begin{boxexo}[title={\ding{45}\quad Exercice~15.9. Demande \niveauexpert}]
Rédigez la demande de preuve adressée au service d'hébergement.
\end{boxexo}

\begin{boxexo}[title={\ding{45}\quad Exercice~15.10. Bilan domaines \niveauexpert}]
Récapitulez ce que chaque domaine (disque, mémoire, réseau, mobile, nuage) a
apporté à TSANAGA.
\end{boxexo}

\bigskip

\begin{boxinfo}
\textbf{Pour aller plus loin.}
\begin{itemize}
    \item ISO/IEC 27037~\cite{iso27037}, 27042~\cite{iso27042}~;
          Budapest~\cite{budapest2001}.
\end{itemize}
\end{boxinfo}

\begin{boxinfo}
\textbf{Transition.} Un suspect averti efface ses traces~: la Partie~V traite
des défis, à commencer par l'anti-forensique
(Chapitre~\ref{chap:anti}).
\end{boxinfo}
}{}

    \part{Les Défis de l'Enquête}
    \label{part:defis}

        \IfFileExists{partie5/P05_C16_antiforensique.tex}{% ============================================================================
% Fichier  : partie5/P05_C16_antiforensique.tex : FORN (LPI342)
% Sources : ISO/IEC 27042 ; NIST SP 800-86. Opération TSANAGA.
% ============================================================================

\chapter{L'Anti-Forensique et le Chiffrement}
\label{chap:anti}
\chaptermark{L'anti-forensique}

\epigraph{%
    « Effacer une trace en laisse souvent une autre~:\\
    celle de l'effacement. »%
}{Principe du cours}

\niveauChapitre{Expert}{Partie~IV.}

\begin{boxobjectifs}
\begin{itemize}
    \item Reconnaître les techniques d'anti-forensique.
    \item Comprendre l'effacement sécurisé et ses limites.
    \item Traiter le chiffrement dans une enquête.
    \item Repérer les traces d'une dissimulation.
    \item Rester honnête sur ce que l'on ne peut établir (NEK).
\end{itemize}
\end{boxobjectifs}

\bigskip

% ============================================================================
\section*{Question directrice : si M. NDONGO a effacé et chiffré ses traces, l'enquête est-elle bloquée ?}
% ============================================================================

Un suspect averti efface, chiffre, dissimule. Ces techniques compliquent
l'enquête mais laissent elles-mêmes des traces. Ce chapitre apprend à les
reconnaître et à dire honnêtement ce qui reste hors d'atteinte.

\medskip
\textbf{Plan du chapitre.}
\begin{enumerate}
    \item Les techniques d'anti-forensique (section~\ref{sec:c16-techniques}).
    \item L'effacement sécurisé (section~\ref{sec:c16-effacement}).
    \item Le chiffrement (section~\ref{sec:c16-chiffrement}).
    \item L'honnêteté du NEK (section~\ref{sec:c16-nek}).
\end{enumerate}

\separateur

% ============================================================================
\section{Les Techniques d'Anti-Forensique}
\label{sec:c16-techniques}
% ============================================================================

\begin{tableaucours}{L{3.6cm}L{8.6cm}}{%
    \tableauHeader{Technique} &
    \tableauHeader{Trace qu'elle laisse souvent}
}
Effacement sécurisé & Absence anormale, outil installé, date d'exécution \\
Chiffrement & Conteneur, extension, outil, clé en mémoire \\
Modification d'horodatage & Incohérences entre artefacts \\
Dissimulation (renommage) & Signature qui trahit le vrai type \\
Effacement des journaux & Événement de nettoyage, trou dans la suite \\
\end{tableaucours}
\vspace{-0.8em}
\begin{center}{\small\itshape Tableau~16.1. Techniques d'anti-forensique et leurs traces}\end{center}

% ============================================================================
\section{L'Effacement Sécurisé}
\label{sec:c16-effacement}
% ============================================================================

L'\termecle{effacement sécurisé} réécrit l'espace pour empêcher la
récupération. Mais il laisse des indices~: l'outil utilisé, sa date
d'exécution, l'absence anormale de données là où il devrait y en avoir. Un
effacement ciblé, juste avant une saisie, est en soi un élément.

% ============================================================================
\section{Le Chiffrement}
\label{sec:c16-chiffrement}
% ============================================================================

\begin{boxalert}
\textbf{Violation typique.} Tenter de « casser » un chiffrement fort par force
brute, sans espoir réaliste.

\smallskip
\textbf{Règle~:} contre un chiffrement fort, on cherche la clé ailleurs
(mémoire capturée à chaud, mot de passe noté, phrase de passe faible), pas la
force brute. Dans TSANAGA, la clé était en mémoire
(Chapitre~\ref{chap:memoire}).
\end{boxalert}

% ============================================================================
\section{L'Honnêteté du NEK}
\label{sec:c16-nek}
% ============================================================================

\begin{boxscenario}{Opération TSANAGA. Ce qui reste hors d'atteinte}
\textbf{Constatation.} Un second conteneur chiffré, non monté à la saisie,
reste inaccessible~: aucune clé trouvée, chiffrement fort.

\medskip
\textbf{Analyse} (confiance élevée sur l'impossibilité). On ne peut établir le
contenu de ce conteneur~: c'est un \sigle{NEK}, à déclarer tel quel. On ne
suppose pas son contenu~; on rapporte qu'il existe et reste chiffré. L'honnêteté
sur les limites renforce la crédibilité du reste.
\end{boxscenario}

\begin{boxCRO}
\textbf{Déclarer l'inaccessible.} Prétendre connaître le contenu d'un
conteneur non ouvert dégraderait la \cro{0}{-}{-}~: affirmation non fiable,
non opposable. Le NEK préserve la valeur de l'ensemble du rapport.
\end{boxCRO}

\begin{boxmetier}{Expert judiciaire}
\textbf{Dire ce qu'on ne sait pas.} L'expert qui déclare honnêtement un NEK
(« ce conteneur reste chiffré, son contenu est indéterminé ») est plus crédible
que celui qui comble les vides par des suppositions. La limite assumée est une
force.
\end{boxmetier}

\separateur

\begin{boxresume}
\textbf{Ce qu'il faut retenir de ce chapitre~:}
\begin{itemize}
    \item \textbf{L'anti-forensique} laisse souvent la trace de
          l'effacement.
    \item \textbf{L'effacement sécurisé} réécrit, mais laisse des indices.
    \item \textbf{Contre le chiffrement fort}, on cherche la clé ailleurs.
    \item \textbf{Ce qui reste inaccessible} est un NEK, à déclarer.
    \item \textbf{L'honnêteté sur les limites} renforce la crédibilité.
\end{itemize}
\end{boxresume}

\begin{boxinfo}
\textbf{Points de contrôle de l'élève.} Je sais~: reconnaître les techniques
d'anti-forensique~; expliquer l'effacement sécurisé~; traiter le
chiffrement~; déclarer un NEK~; défendre l'honnêteté des limites.
\end{boxinfo}

\bigskip

\begin{boxexo}[title={\ding{45}\quad Exercice~16.1. Techniques \niveaubase}]
Citez trois techniques d'anti-forensique.
\end{boxexo}

\begin{boxexo}[title={\ding{45}\quad Exercice~16.2. Effacement \niveaubase}]
Quelles traces laisse un effacement sécurisé~?
\end{boxexo}

\begin{boxexo}[title={\ding{45}\quad Exercice~16.3. Chiffrement \niveaubase}]
Contre un chiffrement fort, où cherche-t-on la clé~?
\end{boxexo}

\begin{boxexo}[title={\ding{45}\quad Exercice~16.4. L'effacement comme indice \niveaumoyen}]
Pourquoi un effacement juste avant une saisie est-il un élément~?
\end{boxexo}

\begin{boxexo}[title={\ding{45}\quad Exercice~16.5. NEK \niveaumoyen}]
Que faire du second conteneur chiffré inaccessible~?
\end{boxexo}

\begin{boxexo}[title={\ding{45}\quad Exercice~16.6. Force brute \niveaumoyen}]
Pourquoi la force brute contre un chiffrement fort est-elle vaine~?
\end{boxexo}

\begin{boxexo}[title={\ding{45}\quad Exercice~16.7. CRO \niveaumoyen}]
Pourquoi supposer le contenu d'un conteneur dégrade-t-il la preuve~?
\end{boxexo}

\begin{boxexo}[title={\ding{45}\quad Exercice~16.8. Honnêteté \niveaumoyen}]
En quoi déclarer un NEK renforce-t-il la crédibilité~?
\end{boxexo}

\begin{boxexo}[title={\ding{45}\quad Exercice~16.9. Détecter \niveauexpert}]
Comment détecter une modification d'horodatage dans TSANAGA~?
\end{boxexo}

\begin{boxexo}[title={\ding{45}\quad Exercice~16.10. Note de limites \niveauexpert}]
Rédigez la section « limites de l'expertise » du rapport TSANAGA.
\end{boxexo}

\bigskip

\begin{boxinfo}
\textbf{Pour aller plus loin.}
\begin{itemize}
    \item ISO/IEC 27042~\cite{iso27042}~; NIST SP~800-86~\cite{nist80086}.
\end{itemize}
\end{boxinfo}

\begin{boxinfo}
\textbf{Transition.} Il faut maintenant assembler les domaines en une
histoire~: le Chapitre~\ref{chap:reconstitution} traite de la corrélation et
de la reconstitution chronologique.
\end{boxinfo}
}{}
        \IfFileExists{partie5/P05_C17_reconstitution.tex}{% ============================================================================
% Fichier  : partie5/P05_C17_reconstitution.tex : FORN (LPI342)
% Sources : ISO/IEC 27042 ; NIST SP 800-86. Opération TSANAGA.
% ============================================================================

\chapter{Corrélation et Reconstitution Chronologique}
\label{chap:reconstitution}
\chaptermark{Reconstitution chronologique}

\epigraph{%
    « Une preuve isolée se discute~;\\
    un faisceau convergent s'impose. »%
}{Principe du cours}

\niveauChapitre{Expert}{Toute la Partie~IV~; Chapitre~\ref{chap:anti}.}

\begin{boxobjectifs}
\begin{itemize}
    \item Construire une chronologie à partir de plusieurs domaines.
    \item Recaler les horloges et traiter les incohérences.
    \item Faire converger les indices en un faisceau.
    \item Distinguer le certain, le probable, l'indéterminé.
    \item Préparer la démonstration pour le rapport.
\end{itemize}
\end{boxobjectifs}

\bigskip

% ============================================================================
\section*{Question directrice : que s'est-il exactement passé, et avec quelle certitude ?}
% ============================================================================

Disque, mémoire, réseau, mobile, nuage ont chacun livré des éléments. Les
assembler en une chronologie cohérente, en pesant la certitude de chacun, est
le cœur de la reconstitution. Ce chapitre l'enseigne.

\medskip
\textbf{Plan du chapitre.}
\begin{enumerate}
    \item Construire la chronologie (section~\ref{sec:c17-chrono}).
    \item Recaler et traiter les incohérences (section~\ref{sec:c17-recaler}).
    \item Le faisceau convergent (section~\ref{sec:c17-faisceau}).
    \item Certain, probable, indéterminé (section~\ref{sec:c17-certitude}).
\end{enumerate}

\separateur

% ============================================================================
\section{Construire la Chronologie}
\label{sec:c17-chrono}
% ============================================================================

\begin{boxprocedure}
\textbf{Construire une chronologie multi-domaines}
\begin{enumerate}
    \item Rassembler les événements datés de tous les domaines.
    \item Recaler les horloges (écarts connus).
    \item Ordonner et repérer les enchaînements.
    \item Noter la source et la certitude de chaque événement.
    \item Dater chaque étape de l'affaire.
\end{enumerate}
\end{boxprocedure}

% ============================================================================
\section{Recaler et Traiter les Incohérences}
\label{sec:c17-recaler}
% ============================================================================

Les horloges diffèrent (fuseau, dérive, modification volontaire). On les
recale sur une référence, on documente les écarts, et l'on traite les
incohérences~: une date impossible (un fichier modifié après son effacement)
signale souvent une manipulation (Chapitre~\ref{chap:anti}).

% ============================================================================
\section{Le Faisceau Convergent}
\label{sec:c17-faisceau}
% ============================================================================

\begin{boxscenario}{Opération TSANAGA. Le faisceau}
\textbf{Constatation.} La chronologie assemble~: 3~janvier, ouverture et copie
de \texttt{clients.xlsx} vers un amovible (disque)~; 5~janvier, téléversement
vers un service d'hébergement (réseau, trace de synchronisation)~; janvier,
messages de négociation (mobile)~; à la saisie, le fichier dans un conteneur
chiffré ouvert (mémoire).

\medskip
\textbf{Analyse} (confiance élevée). Quatre domaines convergent sur la même
histoire~: copie, envoi, négociation, détention. Aucun indice isolé ne
suffirait~; ensemble, ils forment un faisceau solide. Un second conteneur
reste NEK.
\end{boxscenario}

% ============================================================================
\section{Certain, Probable, Indéterminé}
\label{sec:c17-certitude}
% ============================================================================

\begin{boxCRO}
\textbf{Peser chaque élément.} La reconstitution distingue le \emph{certain}
(établi par plusieurs sources), le \emph{probable} (une source, à corroborer)
et l'\emph{indéterminé} (NEK). Attribuer la bonne certitude à chaque
affirmation sert la \cro{0}{+}{+}~: un rapport qui surévalue sa certitude
perd en opposabilité.
\end{boxCRO}

\begin{boxmetier}{Analyste forensique et expert judiciaire}
\textbf{Assembler sans forcer.} L'expert fait converger les domaines, mais ne
force pas le faisceau~: il laisse les NEK comme tels. Un faisceau honnête, qui
assume ses zones d'ombre, convainc davantage qu'une démonstration trop lisse.
\end{boxmetier}

\separateur

\begin{boxresume}
\textbf{Ce qu'il faut retenir de ce chapitre~:}
\begin{itemize}
    \item \textbf{La chronologie} assemble les événements datés des domaines.
    \item \textbf{On recale les horloges} et on traite les incohérences.
    \item \textbf{Le faisceau convergent} fait la force de la preuve.
    \item \textbf{On pèse} certain, probable, indéterminé.
    \item \textbf{On n'efface pas les NEK} pour lisser la démonstration.
\end{itemize}
\end{boxresume}

\begin{boxinfo}
\textbf{Points de contrôle de l'élève.} Je sais~: construire une chronologie
multi-domaines~; recaler les horloges~; bâtir un faisceau~; peser la
certitude~; préserver les NEK.
\end{boxinfo}

\bigskip

\begin{boxexo}[title={\ding{45}\quad Exercice~17.1. Chronologie \niveaubase}]
Quelles étapes pour construire une chronologie multi-domaines~?
\end{boxexo}

\begin{boxexo}[title={\ding{45}\quad Exercice~17.2. Recaler \niveaubase}]
Pourquoi recaler les horloges~?
\end{boxexo}

\begin{boxexo}[title={\ding{45}\quad Exercice~17.3. Faisceau \niveaubase}]
Pourquoi un indice isolé est-il plus fragile qu'un faisceau~?
\end{boxexo}

\begin{boxexo}[title={\ding{45}\quad Exercice~17.4. Incohérence \niveaumoyen}]
Que signale un fichier « modifié après son effacement »~?
\end{boxexo}

\begin{boxexo}[title={\ding{45}\quad Exercice~17.5. TSANAGA \niveaumoyen}]
Dressez le faisceau convergent de l'affaire, domaine par domaine.
\end{boxexo}

\begin{boxexo}[title={\ding{45}\quad Exercice~17.6. Certitude \niveaumoyen}]
Classez en certain, probable, indéterminé trois éléments de l'affaire.
\end{boxexo}

\begin{boxexo}[title={\ding{45}\quad Exercice~17.7. Surévaluer \niveaumoyen}]
Pourquoi surévaluer sa certitude nuit-il au rapport~?
\end{boxexo}

\begin{boxexo}[title={\ding{45}\quad Exercice~17.8. NEK \niveaumoyen}]
Pourquoi ne pas effacer les NEK de la reconstitution~?
\end{boxexo}

\begin{boxexo}[title={\ding{45}\quad Exercice~17.9. Chronologie complète \niveauexpert}]
Rédigez la chronologie complète et datée de l'opération TSANAGA.
\end{boxexo}

\begin{boxexo}[title={\ding{45}\quad Exercice~17.10. Carte de certitude \niveauexpert}]
Construisez la carte « élément / source / certitude » de l'affaire.
\end{boxexo}

\bigskip

\begin{boxinfo}
\textbf{Pour aller plus loin.}
\begin{itemize}
    \item ISO/IEC 27042~\cite{iso27042}~; NIST SP~800-86~\cite{nist80086}.
\end{itemize}
\end{boxinfo}

\begin{boxinfo}
\textbf{Transition.} Le faisceau établi, il faut le restituer~: la Partie~VI
s'ouvre sur le rapport forensique (Chapitre~\ref{chap:rapport}).
\end{boxinfo}
}{}

    \part{Restituer l'Enquête}
    \label{part:restituer}

        \IfFileExists{partie6/P06_C18_rapport.tex}{% ============================================================================
% Fichier  : partie6/P06_C18_rapport.tex : FORN (LPI342)
% Sources : ISO/IEC 27042 ; ACPO. Opération TSANAGA.
% ============================================================================

\chapter{Le Rapport Forensique}
\label{chap:rapport}
\chaptermark{Le rapport forensique}

\epigraph{%
    « Le rapport est la seule partie de l'enquête que le juge verra. »%
}{Principe du cours}

\niveauChapitre{Intermédiaire}{Toute l'enquête~; Chapitre~\ref{chap:reconstitution}.}

\begin{boxobjectifs}
\begin{itemize}
    \item Structurer un rapport forensique clair et complet.
    \item Séparer méthode, constatation et interprétation.
    \item Écrire pour un lecteur non technique sans perdre la rigueur.
    \item Documenter la reproductibilité.
    \item Assumer les limites et les NEK.
\end{itemize}
\end{boxobjectifs}

\bigskip

% ============================================================================
\section*{Question directrice : comment transformer des mois d'enquête en un document que le juge comprendra et pourra opposer ?}
% ============================================================================

L'enquête ne vaut que par son rapport. Il doit être compris par un
magistrat non technicien, et résister à la contestation d'un expert adverse.
Ce chapitre le construit.

\medskip
\textbf{Plan du chapitre.}
\begin{enumerate}
    \item La structure (section~\ref{sec:c18-structure}).
    \item Méthode, fait, interprétation (section~\ref{sec:c18-separer}).
    \item Écrire pour être compris (section~\ref{sec:c18-ecrire}).
    \item Reproductibilité et limites (section~\ref{sec:c18-limites}).
\end{enumerate}

\separateur

% ============================================================================
\section{La Structure}
\label{sec:c18-structure}
% ============================================================================

\begin{tableaucours}{L{4.0cm}L{8.2cm}}{%
    \tableauHeader{Partie} &
    \tableauHeader{Contenu}
}
Synthèse & Ce qui est établi, en termes simples, pour le magistrat \\
Cadre et mission & Réquisition, périmètre, pièces, dates \\
Méthode & Normes suivies, outils, actes, condensats \\
Constatations & Les faits établis, par domaine \\
Interprétation & Analyse, niveaux de confiance, hypothèses \\
Limites et NEK & Ce qui n'a pu être établi \\
Annexes & Chaîne de possession, condensats, journaux \\
\end{tableaucours}
\vspace{-0.8em}
\begin{center}{\small\itshape Tableau~18.1. Structure d'un rapport forensique}\end{center}

% ============================================================================
\section{Méthode, Fait, Interprétation}
\label{sec:c18-separer}
% ============================================================================

\begin{boxalert}
\textbf{Violation typique.} Mêler le fait et l'opinion (« le fichier a été
copié, donc NDONGO est un voleur »).

\smallskip
\textbf{Règle~:} on sépare nettement la \emph{méthode} (comment on a procédé),
la \emph{constatation} (le fait établi) et l'\emph{interprétation} (ce qu'on
en déduit, avec un niveau de confiance). Le juge doit pouvoir distinguer les
trois d'un coup d'œil.
\end{boxalert}

% ============================================================================
\section{Écrire pour être Compris}
\label{sec:c18-ecrire}
% ============================================================================

Le rapport s'adresse d'abord à un non-technicien. On explique les termes, on
illustre, on va du simple au détaillé~: une synthèse accessible en tête, le
détail technique ensuite, les preuves en annexe. La clarté n'est pas l'ennemie
de la rigueur~: elle la sert.

% ============================================================================
\section{Reproductibilité et Limites}
\label{sec:c18-limites}
% ============================================================================

\begin{boxscenario}{Opération TSANAGA. Le rapport remis}
\textbf{Constatation.} Le rapport de M.~BIKORO établit la copie (disque),
l'envoi (réseau et nuage), la négociation (mobile), et déclare le second
conteneur comme NEK. Chaque constat cite sa méthode et son condensat~; la
chaîne de possession est en annexe.

\medskip
\textbf{Analyse} (confiance élevée). Un expert adverse pourrait reproduire
l'analyse et obtenir le même résultat~: c'est la marque d'un rapport
opposable. Les limites assumées (le second conteneur) renforcent sa
crédibilité.
\end{boxscenario}

\begin{boxmetier}{Expert judiciaire}
\textbf{Écrire pour le contradictoire.} L'expert rédige en imaginant un
contradicteur~: chaque affirmation doit pouvoir être vérifiée, chaque méthode
reproduite. Un rapport qui anticipe la contestation la désamorce.
\end{boxmetier}

\separateur

\begin{boxresume}
\textbf{Ce qu'il faut retenir de ce chapitre~:}
\begin{itemize}
    \item \textbf{Le rapport} est la seule partie que le juge verra.
    \item \textbf{On sépare} méthode, constatation, interprétation.
    \item \textbf{On écrit} pour un non-technicien, sans perdre la rigueur.
    \item \textbf{On documente} la reproductibilité.
    \item \textbf{On assume} les limites et les NEK.
\end{itemize}
\end{boxresume}

\begin{boxinfo}
\textbf{Points de contrôle de l'élève.} Je sais~: structurer un rapport~;
séparer méthode, fait, interprétation~; écrire clair~; documenter la
reproductibilité~; assumer les limites.
\end{boxinfo}

\bigskip

\begin{boxexo}[title={\ding{45}\quad Exercice~18.1. Structure \niveaubase}]
Citez les parties d'un rapport forensique.
\end{boxexo}

\begin{boxexo}[title={\ding{45}\quad Exercice~18.2. Trois registres \niveaubase}]
Distinguez méthode, constatation, interprétation.
\end{boxexo}

\begin{boxexo}[title={\ding{45}\quad Exercice~18.3. Lecteur \niveaubase}]
À qui s'adresse d'abord le rapport~?
\end{boxexo}

\begin{boxexo}[title={\ding{45}\quad Exercice~18.4. Mêler \niveaumoyen}]
Pourquoi ne pas mêler le fait et l'opinion~?
\end{boxexo}

\begin{boxexo}[title={\ding{45}\quad Exercice~18.5. Reproductible \niveaumoyen}]
Comment rendre un rapport reproductible~?
\end{boxexo}

\begin{boxexo}[title={\ding{45}\quad Exercice~18.6. NEK \niveaumoyen}]
Pourquoi déclarer le second conteneur comme NEK dans le rapport~?
\end{boxexo}

\begin{boxexo}[title={\ding{45}\quad Exercice~18.7. Synthèse \niveaumoyen}]
Rédigez la synthèse (pour le magistrat) de l'affaire TSANAGA.
\end{boxexo}

\begin{boxexo}[title={\ding{45}\quad Exercice~18.8. Contradictoire \niveaumoyen}]
Comment écrire en anticipant la contestation~?
\end{boxexo}

\begin{boxexo}[title={\ding{45}\quad Exercice~18.9. Rapport complet \niveauexpert}]
Rédigez le plan détaillé du rapport TSANAGA, partie par partie.
\end{boxexo}

\begin{boxexo}[title={\ding{45}\quad Exercice~18.10. Constatation \niveauexpert}]
Rédigez une constatation complète (copie du fichier), méthode et condensat
inclus.
\end{boxexo}

\bigskip

\begin{boxinfo}
\textbf{Pour aller plus loin.}
\begin{itemize}
    \item ISO/IEC 27042~\cite{iso27042}~; ACPO~\cite{acpo2012}.
\end{itemize}
\end{boxinfo}

\begin{boxinfo}
\textbf{Transition.} Le rapport se défend à l'oral~: le
Chapitre~\ref{chap:tribunal} traite de la présentation devant un tribunal.
\end{boxinfo}
}{}
        \IfFileExists{partie6/P06_C19_tribunal.tex}{% ============================================================================
% Fichier  : partie6/P06_C19_tribunal.tex : FORN (LPI342)
% Sources : ACPO ; déontologie de l'expertise. Opération TSANAGA.
% ============================================================================

\chapter{La Présentation devant un Tribunal}
\label{chap:tribunal}
\chaptermark{Devant un tribunal}

\epigraph{%
    « Devant le juge, l'expert ne plaide pas~: il explique. »%
}{Principe du cours}

\niveauChapitre{Intermédiaire à expert}{Chapitre~\ref{chap:rapport}.}

\begin{boxobjectifs}
\begin{itemize}
    \item Préparer une présentation orale claire.
    \item Expliquer une preuve technique à un public non technicien.
    \item Répondre aux questions sans sortir de son rôle.
    \item Tenir face à la contradiction.
    \item Rester dans les limites de l'expertise.
\end{itemize}
\end{boxobjectifs}

\bigskip

% ============================================================================
\section*{Question directrice : comment faire comprendre au tribunal ce que montrent les preuves, sans les trahir ?}
% ============================================================================

Le rapport est écrit~; il faut maintenant l'exposer à l'audience, répondre aux
questions, tenir face à un expert adverse. Ce chapitre prépare cette
restitution orale.

\medskip
\textbf{Plan du chapitre.}
\begin{enumerate}
    \item Préparer la présentation (section~\ref{sec:c19-preparer}).
    \item Expliquer sans trahir (section~\ref{sec:c19-expliquer}).
    \item Répondre et tenir (section~\ref{sec:c19-repondre}).
    \item Rester dans son rôle (section~\ref{sec:c19-role}).
\end{enumerate}

\separateur

% ============================================================================
\section{Préparer la Présentation}
\label{sec:c19-preparer}
% ============================================================================

\begin{boxprocedure}
\textbf{Préparer une présentation au tribunal}
\begin{enumerate}
    \item Résumer les faits établis en quelques phrases simples.
    \item Préparer une explication imagée de chaque notion technique.
    \item Anticiper les questions et les objections.
    \item Préparer les annexes (chaîne de possession, condensats).
    \item Distinguer, à l'oral aussi, le certain, le probable, le NEK.
\end{enumerate}
\end{boxprocedure}

% ============================================================================
\section{Expliquer sans Trahir}
\label{sec:c19-expliquer}
% ============================================================================

Expliquer un condensat, une image bit à bit ou une chaîne de possession à un
magistrat demande des images justes, sans simplification trompeuse. On dit
« une empreinte qui change entièrement si un seul caractère change », non « un
code infaillible ». La pédagogie ne doit jamais déformer le fait.

% ============================================================================
\section{Répondre et Tenir}
\label{sec:c19-repondre}
% ============================================================================

\begin{boxscenario}{Opération TSANAGA. La question de la défense}
\textbf{Constatation.} À l'audience, l'avocat de la défense demande~:
« Comment être sûr que le disque analysé est bien celui de mon client~? »

\medskip
\textbf{Analyse} (confiance élevée). M.~BIKORO répond par la chaîne de
possession et le condensat~: le disque scellé, la fiche de scellé, l'empreinte
identique à l'acquisition et à l'audience. La réponse est factuelle, vérifiable,
calme. Il reconnaît aussi, sans détour, la fragilité de la clé USB et le NEK
du second conteneur.
\end{boxscenario}

% ============================================================================
\section{Rester dans son Rôle}
\label{sec:c19-role}
% ============================================================================

\begin{boxalert}
\textbf{Violation typique.} Conclure à la culpabilité du prévenu.

\smallskip
\textbf{Règle~:} l'expert établit des faits techniques et leur portée~; il ne
qualifie pas l'infraction ni la culpabilité, qui relèvent du juge. Sortir de
ce rôle décrédibilise l'expertise.
\end{boxalert}

\begin{boxmetier}{Expert judiciaire}
\textbf{Calme, clair, dans son rôle.} À la barre, l'expert reste factuel, ne
se laisse pas entraîner dans le débat juridique, reconnaît les limites, et
renvoie toujours à ses pièces. La sérénité et l'honnêteté sont ses meilleures
défenses.
\end{boxmetier}

\separateur

\begin{boxresume}
\textbf{Ce qu'il faut retenir de ce chapitre~:}
\begin{itemize}
    \item \textbf{On prépare} une synthèse simple et des explications imagées.
    \item \textbf{On explique sans trahir} le fait.
    \item \textbf{On répond} par les pièces (chaîne, condensat), calmement.
    \item \textbf{On reconnaît} les limites et les NEK.
    \item \textbf{L'expert ne conclut pas} à la culpabilité.
\end{itemize}
\end{boxresume}

\begin{boxinfo}
\textbf{Points de contrôle de l'élève.} Je sais~: préparer une présentation~;
expliquer une notion technique~; répondre par les pièces~; tenir face à la
contradiction~; rester dans mon rôle.
\end{boxinfo}

\bigskip

\begin{boxexo}[title={\ding{45}\quad Exercice~19.1. Préparer \niveaubase}]
Citez trois éléments à préparer avant une audience.
\end{boxexo}

\begin{boxexo}[title={\ding{45}\quad Exercice~19.2. Expliquer \niveaubase}]
Comment expliquer un condensat à un magistrat~?
\end{boxexo}

\begin{boxexo}[title={\ding{45}\quad Exercice~19.3. Le rôle \niveaubase}]
L'expert peut-il conclure à la culpabilité~?
\end{boxexo}

\begin{boxexo}[title={\ding{45}\quad Exercice~19.4. La question \niveaumoyen}]
Répondez à « comment être sûr que c'est bien le disque de mon client~? ».
\end{boxexo}

\begin{boxexo}[title={\ding{45}\quad Exercice~19.5. Simplifier \niveaumoyen}]
Pourquoi une simplification trompeuse est-elle dangereuse~?
\end{boxexo}

\begin{boxexo}[title={\ding{45}\quad Exercice~19.6. Limites \niveaumoyen}]
Pourquoi reconnaître à l'oral la fragilité de la clé USB~?
\end{boxexo}

\begin{boxexo}[title={\ding{45}\quad Exercice~19.7. Contradiction \niveaumoyen}]
Comment tenir face à un expert adverse~?
\end{boxexo}

\begin{boxexo}[title={\ding{45}\quad Exercice~19.8. Sérénité \niveaumoyen}]
Pourquoi le calme est-il une force à la barre~?
\end{boxexo}

\begin{boxexo}[title={\ding{45}\quad Exercice~19.9. Script \niveauexpert}]
Préparez le déroulé de la présentation orale de TSANAGA.
\end{boxexo}

\begin{boxexo}[title={\ding{45}\quad Exercice~19.10. Objections \niveauexpert}]
Listez cinq objections probables et vos réponses factuelles.
\end{boxexo}

\bigskip

\begin{boxinfo}
\textbf{Pour aller plus loin.}
\begin{itemize}
    \item ACPO~\cite{acpo2012}~; ISO/IEC 27042~\cite{iso27042}.
\end{itemize}
\end{boxinfo}

\begin{boxinfo}
\textbf{Transition.} Reste à progresser~: le Chapitre~\ref{chap:mesurer} clôt
le cours sur la mesure de la qualité de l'enquête.
\end{boxinfo}
}{}
        \IfFileExists{partie6/P06_C20_mesurer.tex}{% ============================================================================
% Fichier  : partie6/P06_C20_mesurer.tex : FORN (LPI342)
% Sources : ISO/IEC 27043 ; cadre OUI/NON/NEK (Minka). Opération TSANAGA.
% ============================================================================

\chapter{Mesurer et Progresser}
\label{chap:mesurer}
\chaptermark{Mesurer et conclure}

\epigraph{%
    « La qualité d'une enquête se mesure aussi\\
    à l'honnêteté de ses incertitudes. »%
}{Principe du cours}

\niveauChapitre{Avancé}{L'ensemble du cours~; notions d'auditologie.}

\begin{boxobjectifs}
\begin{itemize}
    \item Définir la qualité d'une enquête forensique.
    \item Choisir des indicateurs honnêtes.
    \item Appliquer le verdict OUI, NON, NEK.
    \item Conduire une revue par les pairs.
    \item Ouvrir sur la recherche.
\end{itemize}
\end{boxobjectifs}

\bigskip

% ============================================================================
\section*{Question directrice : l'enquête TSANAGA était-elle « bonne » ?}
% ============================================================================

Une enquête se juge non à sa conclusion, mais à sa méthode, sa traçabilité et
l'honnêteté de ses limites. Ce chapitre clôt le cours sur la mesure de la
qualité.

\medskip
\textbf{Plan du chapitre.}
\begin{enumerate}
    \item La qualité d'une enquête (section~\ref{sec:c20-qualite}).
    \item Des indicateurs honnêtes (section~\ref{sec:c20-indicateurs}).
    \item Le verdict OUI / NON / NEK (section~\ref{sec:c20-verdict}).
    \item Revue par les pairs et recherche (section~\ref{sec:c20-pairs}).
\end{enumerate}

\separateur

% ============================================================================
\section{La Qualité d'une Enquête}
\label{sec:c20-qualite}
% ============================================================================

\begin{boxdef}
\termecle{Qualité d'une enquête forensique}~: aptitude de l'enquête à produire
des conclusions \emph{fiables} et \emph{opposables}, par une méthode tracée,
reproductible, et une honnêteté assumée sur les limites. Une « bonne »
conclusion issue d'une mauvaise méthode ne vaut rien.
\end{boxdef}

% ============================================================================
\section{Des Indicateurs Honnêtes}
\label{sec:c20-indicateurs}
% ============================================================================

\begin{tableaucours}{L{4.8cm}L{7.4cm}}{%
    \tableauHeader{Indicateur} &
    \tableauHeader{Ce qu'il dit}
}
Traçabilité des actes & Chaque geste documenté et reproductible \\
Intégrité des pièces & Condensats vérifiés à chaque étape \\
Part des conclusions corroborées & Faits tenus par plusieurs sources \\
Clarté des limites & NEK déclarés, incertitudes assumées \\
Tenue en contradictoire & Pièces résistant à la contestation \\
\end{tableaucours}
\vspace{-0.8em}
\begin{center}{\small\itshape Tableau~20.1. Indicateurs de qualité d'une enquête}\end{center}

\begin{boxalert}
\textbf{Indicateur trompeur.} « L'enquête a conclu à la culpabilité. »

\smallskip
\textbf{Règle~:} la conclusion n'est pas un indicateur de qualité~; l'expert
ne conclut d'ailleurs pas à la culpabilité. Ce qui se mesure, c'est la méthode
et l'honnêteté, non le résultat judiciaire.
\end{boxalert}

% ============================================================================
\section{Le Verdict OUI / NON / NEK}
\label{sec:c20-verdict}
% ============================================================================

\begin{boxdef}
\termecle{Verdict à trois valeurs}~: pour chaque question d'enquête, la réponse
est \sigle{OUI} (établi, avec preuve), \sigle{NON} (réfuté, avec preuve) ou
\sigle{NEK} (indéterminé, faute de preuves). Le NEK n'affaiblit pas l'enquête~:
il en garantit l'honnêteté.
\end{boxdef}

\begin{tableaucours}{L{2.0cm}L{4.8cm}L{5.4cm}}{%
    \tableauHeader{Verdict} &
    \tableauHeader{Question (TSANAGA)} &
    \tableauHeader{Fondement}
}
OUI & Le fichier a-t-il été copié vers un amovible~? & Disque et artefacts \\
OUI & A-t-il été téléversé vers un service~? & Réseau et synchronisation \\
NEK & Que contient le second conteneur chiffré~? & Clé introuvable \\
\end{tableaucours}
\vspace{-0.8em}
\begin{center}{\small\itshape Tableau~20.2. Le verdict OUI / NON / NEK appliqué~\cite{minka2026auditologie}}\end{center}

% ============================================================================
\section{Revue par les Pairs et Recherche}
\label{sec:c20-pairs}
% ============================================================================

\begin{boxscenario}{Opération TSANAGA. La relecture}
\textbf{Constatation.} Avant remise, le rapport est relu par un second expert
du cabinet, qui vérifie méthode, condensats, et la juste attribution des
certitudes.

\medskip
\textbf{Analyse} (confiance élevée). La revue par les pairs détecte les
conclusions trop assurées et les NEK oubliés. Elle est, pour l'enquête
forensique, l'équivalent du contrôle qualité~: elle améliore la fiabilité et
l'opposabilité.
\end{boxscenario}

La recherche prolonge le cours~: automatisation de l'analyse, forensique du
nuage et du mobile en contexte africain, et application du cadre de
l'auditologie à la réduction de l'incertitude dans l'enquête numérique.

\begin{boxmetier}{Expert judiciaire et analyste forensique}
\textbf{Se faire relire.} L'expert ne remet jamais un rapport sans relecture
par un pair~: un regard neuf voit les suppositions déguisées en faits. La
qualité se construit à deux.
\end{boxmetier}

\separateur

\begin{boxresume}
\textbf{Ce qu'il faut retenir de ce chapitre~:}
\begin{itemize}
    \item \textbf{La qualité} tient à la méthode et à l'honnêteté, non à la
          conclusion.
    \item \textbf{Des indicateurs}~: traçabilité, intégrité, corroboration,
          limites.
    \item \textbf{OUI, NON, NEK}~: le verdict honnête.
    \item \textbf{La revue par les pairs} est le contrôle qualité de
          l'enquête.
    \item \textbf{La recherche} prolonge la discipline.
\end{itemize}
\end{boxresume}

\begin{boxinfo}
\textbf{Points de contrôle de l'élève.} Je sais~: définir la qualité d'une
enquête~; choisir des indicateurs honnêtes~; appliquer OUI/NON/NEK~; conduire
une revue par les pairs~; situer la recherche.
\end{boxinfo}

\bigskip

\begin{boxexo}[title={\ding{45}\quad Exercice~20.1. Qualité \niveaubase}]
À quoi tient la qualité d'une enquête forensique~?
\end{boxexo}

\begin{boxexo}[title={\ding{45}\quad Exercice~20.2. Trois valeurs \niveaubase}]
Donnez le sens de OUI, NON et NEK.
\end{boxexo}

\begin{boxexo}[title={\ding{45}\quad Exercice~20.3. Trompeur \niveaubase}]
Pourquoi « l'enquête a conclu à la culpabilité » n'est-il pas un indicateur de
qualité~?
\end{boxexo}

\begin{boxexo}[title={\ding{45}\quad Exercice~20.4. Classer \niveaumoyen}]
Attribuez OUI, NON ou NEK aux trois questions de TSANAGA du tableau.
\end{boxexo}

\begin{boxexo}[title={\ding{45}\quad Exercice~20.5. Pairs \niveaumoyen}]
Que détecte une revue par les pairs~?
\end{boxexo}

\begin{boxexo}[title={\ding{45}\quad Exercice~20.6. Méthode ou résultat \niveaumoyen}]
Pourquoi mesure-t-on la méthode et non le résultat judiciaire~?
\end{boxexo}

\begin{boxexo}[title={\ding{45}\quad Exercice~20.7. Indicateurs \niveaumoyen}]
Choisissez trois indicateurs de qualité pour le cabinet NYONG SÉCURITÉ.
\end{boxexo}

\begin{boxexo}[title={\ding{45}\quad Exercice~20.8. NEK \niveaumoyen}]
Pourquoi le NEK garantit-il l'honnêteté de l'enquête~?
\end{boxexo}

\begin{boxexo}[title={\ding{45}\quad Exercice~20.9. Grille de qualité \niveauexpert}]
Construisez une grille d'évaluation de la qualité d'une enquête forensique.
\end{boxexo}

\begin{boxexo}[title={\ding{45}\quad Exercice~20.10. Question de recherche \niveauexpert}]
Formulez une question de recherche sur la forensique du mobile en contexte
africain.
\end{boxexo}

\bigskip

\begin{boxinfo}
\textbf{Pour aller plus loin.}
\begin{itemize}
    \item ISO/IEC 27043~\cite{iso27043}~; cadre de
          l'auditologie~\cite{minka2026auditologie}.
\end{itemize}
\end{boxinfo}

\begin{boxinfo}
\textbf{Fin du cours.} De la scène saisie au verdict mesuré, FORN aura montré
qu'une enquête numérique établit \emph{et} prouve~: acquérir sans altérer,
analyser chaque domaine, reconstituer un faisceau, et restituer avec rigueur,
sans jamais confondre ce que l'on constate et ce que l'on en déduit.
\end{boxinfo}
}{}

% --------------------------------------------------------------------------
% ANNEXES
% --------------------------------------------------------------------------
\appendix

    %% Annexe A : Glossaire
\chapter{Glossaire}
\label{ann:glossaire}

Ce glossaire reprend les notions définies dans les chapitres. Le renvoi
indique le chapitre où la notion est présentée.

\begin{description}
    \item[Acquisition à chaud] Capture de données sur un système en fonctionnement, avant son extinction (Ch.~\ref{chap:acquisition}).
    \item[Bloqueur en écriture] Dispositif matériel ou logiciel qui autorise la lecture d'un support mais interdit toute modification (Ch.~\ref{chap:preserver}).
    \item[Chaîne de possession] Continuité documentée de la détention d'une pièce, de la saisie à la présentation (Ch.~\ref{chap:custody}).
    \item[Condensat] Empreinte de longueur fixe produite par une fonction de hachage, qui prouve l'intégrité (Ch.~\ref{chap:duplication}).
    \item[Effacement sécurisé] Réécriture de l'espace pour empêcher la récupération, laissant néanmoins des indices (Ch.~\ref{chap:anti}).
    \item[Fonction de hachage] Fonction transformant une donnée en un condensat, sensible au moindre bit (MD5, SHA-256) (Ch.~\ref{chap:duplication}).
    \item[Forensique numérique] Investigation numérique à vocation probatoire, soumise à la rigueur de la preuve (Ch.~\ref{chap:forensique}).
    \item[Image forensique] Copie bit à bit d'un support, reproduisant fichiers, espace non alloué et métadonnées (Ch.~\ref{chap:duplication}).
    \item[Impartialité] Disposition à suivre les preuves où elles mènent, sans favoriser une thèse (Ch.~\ref{chap:ethique}).
    \item[Investigation numérique] Démarche générale de collecte, d'analyse et d'interprétation d'éléments numériques (Ch.~\ref{chap:forensique}).
    \item[Métadonnées] Données portées par un fichier (dates, nom, propriétaire), utiles pour dater une action (Ch.~\ref{chap:disque}).
    \item[Neutralité] Absence d'intérêt personnel de l'enquêteur dans l'issue (Ch.~\ref{chap:ethique}).
    \item[Récupération par signatures] Reconstitution de fichiers par leurs motifs de début et de fin, sans l'index (Ch.~\ref{chap:artefacts}).
    \item[Recevabilité] Aptitude d'une preuve à être admise par la justice, selon la manière dont elle a été obtenue et préservée (Ch.~\ref{chap:recevabilite}).
    \item[Signature (de fichier)] Octets de début d'un fichier qui trahissent son type, indépendamment de l'extension (Ch.~\ref{chap:artefacts}).
    \item[Système de fichiers] Organisation des données d'une partition (FAT, NTFS, EXT) (Ch.~\ref{chap:disque}).
    \item[Trilemme CRO] Confidentialité, Fiabilité, Opposabilité~: les trois dimensions de toute trace et de tout acte d'enquête (Ch.~\ref{chap:prologue}).
    \item[Verdict à trois valeurs] OUI (établi), NON (réfuté), NEK (indéterminé) (Ch.~\ref{chap:mesurer}).
\end{description}

    %% Annexe B : Référentiel des outils
\chapter{Référentiel des Outils}
\label{ann:outils}

Les outils cités sont des exemples. Le cours insiste sur la méthode~: l'outil
sert l'enquête, il ne la remplace pas. Tout usage se fait dans un cadre
autorisé, sur des copies.

\begin{tableaucours}{L{4.2cm}L{5.4cm}L{2.8cm}}{%
    \tableauHeader{Catégorie d'outil} &
    \tableauHeader{Usage} &
    \tableauHeader{Chapitre, lab}
}
Bloqueur en écriture & Lire un support sans l'altérer & Ch.~\ref{chap:preserver} \\
Duplication et hachage & Image bit à bit, condensat (SHA-256) & Ch.~\ref{chap:duplication}, FORN-LAB-02 \\
Examen de disque & Systèmes de fichiers, artefacts & Ch.~\ref{chap:disque}, FORN-LAB-03 \\
Récupération par signatures & Fichiers supprimés, fragments & Ch.~\ref{chap:artefacts}, FORN-LAB-04 \\
Analyse mémoire & Processus, connexions, secrets & Ch.~\ref{chap:memoire}, FORN-LAB-05 \\
Analyse de flux & Journaux de flux exportés & Ch.~\ref{chap:reseau}, FORN-LAB-06 \\
Extraction mobile & Messages, applications, géolocalisation & Ch.~\ref{chap:mobile}, FORN-LAB-07 \\
Tableur & Chronologie, faisceau, qualité & Ch.~\ref{chap:reconstitution}, FORN-LAB-09 \\
\end{tableaucours}

\begin{boxalert}
\textbf{Rappel.} Aucun outil ne s'applique à un support réel sans titre.
L'analyse se fait sur des copies vérifiées par condensat, dans le périmètre de
la réquisition.
\end{boxalert}

    %% Annexe C : Modèles et procédures
\chapter{Modèles et Procédures}
\label{ann:templates}

\section*{Fiche de scellé}
\begin{tableaucours}{L{3.4cm}L{8.8cm}}{%
    \tableauHeader{Rubrique} & \tableauHeader{Contenu}
}
Numéro de scellé & Identifiant unique \\
Pièce & Description, photographie \\
Saisie & Date, lieu, agent \\
Empreinte & Condensat à la réception \\
Transferts & De qui à qui, quand, pourquoi \\
\end{tableaucours}

\section*{Fiche d'acquisition (image)}
\begin{tableaucours}{L{3.4cm}L{8.8cm}}{%
    \tableauHeader{Rubrique} & \tableauHeader{Contenu}
}
Support & Référence, type \\
Méthode & Bloqueur, outil, image bit à bit \\
Empreintes & Condensat original et copie \\
Opérateur & Qui, quand \\
\end{tableaucours}

\section*{Fiche de constat}
\begin{tableaucours}{L{3.4cm}L{8.8cm}}{%
    \tableauHeader{Rubrique} & \tableauHeader{Contenu}
}
Constatation & Le fait établi \\
Méthode & Comment il a été établi, norme suivie \\
Preuve & Artefact, condensat, source \\
Interprétation & Analyse, niveau de confiance \\
\end{tableaucours}

\section*{Plan du rapport forensique}
\begin{enumerate}
    \item Synthèse pour le magistrat.
    \item Cadre et mission (réquisition, périmètre).
    \item Méthode (normes, outils, actes, condensats).
    \item Constatations, par domaine.
    \item Interprétation et niveaux de confiance.
    \item Limites et NEK.
    \item Annexes (chaîne de possession, condensats, journaux).
\end{enumerate}

    %% Annexe D : Textes juridiques et normatifs
\chapter{Textes Juridiques et Normatifs}
\label{ann:textes}

Seuls les éléments vérifiés dans le texte officiel sont cités dans les
chapitres. Ce tableau indique où les trouver.

\section*{Droit camerounais et africain}
\begin{tableaucours}{L{5.0cm}L{4.6cm}L{2.6cm}}{%
    \tableauHeader{Texte} & \tableauHeader{Objet} & \tableauHeader{Chapitre}
}
Loi n° 2010/012 du 21 déc. 2010 & Cybersécurité, cybercriminalité, preuve & \ref{chap:droit} \\
Loi n° 2024/017 du 23 déc. 2024 & Protection des données personnelles & \ref{chap:droit} \\
Convention de Malabo (UA, 2014) & Cybersécurité et données, Afrique & \ref{chap:droit} \\
Convention de Budapest (2001) & Preuve électronique, coopération & \ref{chap:droit} \\
\end{tableaucours}

\section*{Normes et guides internationaux}
\begin{tableaucours}{L{4.4cm}L{5.2cm}L{2.6cm}}{%
    \tableauHeader{Référence} & \tableauHeader{Objet} & \tableauHeader{Chapitre}
}
ISO/IEC 27037 & Acquisition et préservation & \ref{chap:preserver} \\
ISO/IEC 27042 & Analyse et interprétation & \ref{chap:disque} \\
ISO/IEC 27043 & Principes de l'investigation & \ref{chap:demarche} \\
NIST SP 800-86 & Forensique et réponse & \ref{chap:duplication} \\
NIST SP 800-101 Rev. 1 & Forensique mobile & \ref{chap:mobile} \\
RFC 3227 & Collecte et archivage des preuves & \ref{chap:acquisition} \\
ACPO Good Practice Guide & Bonnes pratiques de la preuve numérique & \ref{chap:recevabilite} \\
\end{tableaucours}

    %% Annexe E : Bibliographie commentée
\chapter{Bibliographie Commentée}
\label{ann:bibliographie}

Pour chaque partie, quelques lectures. Les références complètes figurent dans
la bibliographie du livre.

\section*{Partie I · Fondements}
\begin{itemize}
    \item E.~Locard~\cite{locard1920}~: « tout contact laisse une trace ».
    \item ISO/IEC 27043~\cite{iso27043}~: les principes de l'investigation.
    \item ACPO~\cite{acpo2012}~: les bonnes pratiques de la preuve.
\end{itemize}

\section*{Partie II · Cadre juridique et recevabilité}
\begin{itemize}
    \item Lois 2010/012~\cite{loi2010012} et 2024/017~\cite{loi2024017}.
    \item Malabo~\cite{malabo2014}~; Budapest~\cite{budapest2001}.
    \item ISO/IEC 27037~\cite{iso27037}~: la préservation.
\end{itemize}

\section*{Partie III · Acquisition}
\begin{itemize}
    \item NIST SP 800-86~\cite{nist80086}~; RFC 3227~\cite{rfc3227}.
    \item ISO/IEC 27037~\cite{iso27037}~: l'acquisition.
\end{itemize}

\section*{Partie IV · Analyse par domaine}
\begin{itemize}
    \item ISO/IEC 27042~\cite{iso27042}~: l'analyse.
    \item NIST SP 800-101r1~\cite{nist800101}~: la forensique mobile.
    \item MITRE ATT\&CK~\cite{mitreattack}~: les techniques.
\end{itemize}

\section*{Parties V et VI · Défis et restitution}
\begin{itemize}
    \item ISO/IEC 27042~\cite{iso27042}~; ACPO~\cite{acpo2012}.
    \item T.~E.~Minka Mi Nguidjoï~\cite{minka2026auditologie}~: le verdict à trois valeurs.
\end{itemize}

    %% Annexe F : Contacts et ressources
\chapter{Contacts et Ressources}
\label{ann:contacts}

Les adresses sont celles des sites institutionnels à la date de rédaction~; à
vérifier.

\begin{tableaucours}{L{4.4cm}L{4.8cm}L{3.2cm}}{%
    \tableauHeader{Organisation} & \tableauHeader{Rôle} & \tableauHeader{Site}
}
ANTIC & Sécurité électronique, CERT national (Cameroun) & \texttt{www.antic.cm} \\
Union africaine & Convention de Malabo & \texttt{au.int} \\
Conseil de l'Europe & Convention de Budapest & \texttt{coe.int} \\
ISO & Normes ISO et ISO/IEC & \texttt{www.iso.org} \\
NIST & Publications SP 800 & \texttt{www.nist.gov} \\
MITRE & ATT\&CK & \texttt{attack.mitre.org} \\
\end{tableaucours}

\begin{boxinfo}
\textbf{La plateforme du cours.} Les fiches d'auto-évaluation et les dix labs
de NDAP SUKULU pour ce cours sont accessibles par le code \texttt{LPI342}.
\end{boxinfo}


% --------------------------------------------------------------------------
% ÉLÉMENTS FINAUX
% --------------------------------------------------------------------------
\backmatter

    \nocite{*}
    \printbibliography[heading=bibintoc, title={Références bibliographiques}]

    \printindex

% ==============================================================================
% 4ÈME DE COUVERTURE
% ==============================================================================
\clearpage
\thispagestyle{empty}
\pagecolor{CouleurPrimaire}

\begin{tikzpicture}[remember picture, overlay]
    \fill[CouleurAccent]
        (current page.south west) rectangle
        ([yshift=1.8cm] current page.south east);
    \fill[CouleurAccent]
        ([yshift=-1.4cm] current page.north west) rectangle
        ([yshift=-1.7cm] current page.north east);
\end{tikzpicture}

\vspace*{2.5cm}

\begin{center}
    {\color{CouleurAccent}\rule{0.6\linewidth}{1pt}}\\[1.5ex]
    {\color{white}\Large\bfseries À propos de ce cours}\\[1.5ex]
    {\color{CouleurAccent}\rule{0.6\linewidth}{1pt}}
\end{center}

\vspace{1cm}

\begin{center}
\begin{minipage}{0.75\textwidth}
\color{white!85!gray}\setstretch{1.4}
Quand un système d'information s'arrête, la question décisive n'est pas
technique : elle porte sur qui avait décidé, et qui devait rendre compte.
Ce cours apprend à y répondre avant la panne.

\medskip

Il progresse en six parties, précédées d'un prologue : les fondements de la
gouvernance des systèmes d'information, le référentiel COBIT~2019, le cadre
normatif et juridique, ITIL~4 et les processus de gestion des services, les
défis de l'agilité, du nuage et de l'intelligence artificielle, et enfin la
recherche.

\medskip

Chaque notion est ancrée à la fois dans les standards internationaux (ISO,
ISACA, IIA, INTOSAI) et dans le droit camerounais et communautaire (CEMAC,
OHADA). Le cours s'adresse aux étudiants du Master Professionnel en
Management des Systèmes d'Information et aux professionnels qui exercent comme
DSI, responsables GRC, auditeurs des SI, responsables des services
informatiques ou RSSI.
\end{minipage}
\end{center}

\vfill

\begin{center}
    {\color{CouleurAccent}\rule{0.7\linewidth}{1.2pt}}\\[1.2ex]
    {\color{white}\large\bfseries MINKA MI NGUIDJOÏ Thierry Emmanuel}\\[0.4ex]
    {\color{white!80!gray}\small
        CISA~·~CISM~·~CGEIT~·~CRISC~·~CDPSE~·~ISO 27001~·~ISO 22301~·~ISO 9001}\\[0.3ex]
    {\color{white!70!gray}\footnotesize\itshape
        Chercheur au laboratoire d'Ingénierie Mathématiques et des Systèmes d'Information}\\[0.3ex]
    {\color{white!70!gray}\footnotesize\itshape
        Chercheur~|~Cryptologue~|~Auditeur Senior~|~Expert Judiciaire~|~GRC \& Sécurité de l'Information}\\[1.2ex]
    {\color{CouleurAccent}\rule{0.7\linewidth}{1.2pt}}
\end{center}

\vspace{1cm}

\afterpage{\nopagecolor}
\end{document}
